\subsection{\texorpdfstring{关系式 \(\sum_{i} e_i f_i \leqslant n\)}{关系式 Σ ei fi ≤ n}}

设 K 为一个域，v 为 K 上的一个赋值，L 为 K 的一个次数为 n 的有限扩张。~设 \((v_{1}^{\prime}, \ldots, v_{r}^{\prime})\) 为 L 上延拓 v 的赋值，其中任意两个均不等价；若它们独立（当 v 的高为 1 时总是如此），则

\[
\sum_ {i = 1} ^ {r} e (v _ {i} ^ {\prime} / v) f (v _ {i} ^ {\prime} / v) \leqslant n
\]

由命题 2（公式 (6) 和 (7)）可知。我们将看到，这个结果在一般情形下仍然成立。准确地说：

定理 1. 设 K 为一个域，v 为 K 上的一个赋值，L 为 K 的一个次数为 n 的有限扩张。~则：

\begin{enumerate}
\def\labelenumi{(\alph{enumi})}
\item
  每个完备延拓系 \((v_i')_{i\in I}\) 的 \(v\) 到 \(\mathbf{L}\) 的延拓都是有限的。
\item
  \(\sum_{i\in I}e(v_i^{\prime}|v)f(v_i^{\prime}|v)\leqslant n\) 从而 Card(I) \(\leqslant n.\)
\item
  任意两个 \(v_{i}'\) 的环关于包含关系均不可比较。
\end{enumerate}

由于 \(v\) 为非真赋值时定理显然成立，我们设 \(v\) 不为非真赋值。设 \((v_1', \ldots, v_s')\) 为 \(L\) 上延拓 \(v\) 的任意有限族，其中任意两个均不等价。我们首先证明 \(\sum_{i=1}^{s} e(v_i'/v)f(v_i'/v) \leqslant n\)。这将证明 (a) 和 (b)。

我们对 s 归纳，并因此假设不等式已对 0、1、\ldots、s - 1 个赋值的情形成立。分两种情形。

\begin{enumerate}
\def\labelenumi{(\arabic{enumi})}
\tightlist
\item
  假设至少存在两个独立赋值 \(v_i'\)。则存在（§ 7，第 2 节，注 1）一个 \([1, s] = I_1 \cup \cdots \cup I_t\) 对 \([1, s]\) 的划分，使得：
\end{enumerate}

\begin{enumerate}
\def\labelenumi{(\roman{enumi})}
\item
  赋值 \(v_{i}^{\prime}\) 和 \(v_{j}^{\prime}\) 相关的充要条件是 \(i\) 和 \(j\) 属于同一个 \(I_{k}\)；
\item
  \(\operatorname{Card}(\mathbf{I}_k) < s\) 对所有 \(k\) 均成立。
\end{enumerate}

我们在每个 \(I_{k}\) 中选取一个指标 \(i(k)\)。令 \(\hat{L}_{i(k)}\) 表示 L 关于 \(v_{i(k)}^{\prime}\) 的完备化，并令 \(n(k) = [\hat{L}_{i(k)} : \hat{K}]\)。对于所有 \(i \in I_{k}\)，\(v_{i}^{\prime}\) 在 L 上定义出与 \(v_{i(k)}^{\prime}\) 相同的拓扑（§ 7，第 2 节，命题 3），因而可将其延拓为赋值 \(\hat{v}_{i}^{\prime}\)，定义在 \(\hat{L}_{i(k)}\) 上，其在 \(\hat{K}\) 上的限制是 \(\hat{v}\)。由于 \(v_{i}^{\prime}\)（\(i \in I_{k}\)）中任意两个均不等价，\(\hat{v}_{i}^{\prime}\) 也具有同样性质。将归纳假设应用于有序对（\(\hat{K}, \hat{L}_{i(k)}\)），根据命题 2 (a) 的公式 (4)，得到 \(\sum_{i\in I_k}e(v_i^{\prime}|v)f(v_i^{\prime}|v)\leqslant n(k)\)。由于 \(\sum_{k = 1}^{t}n(k)\leqslant n\)（命题 2 (b)，公式

(7)），当然有 \(\sum_{i=1}^{s} e(v_i'/v) f(v_i'/v) \leqslant n\)。

\begin{enumerate}
\def\labelenumi{(\arabic{enumi})}
\setcounter{enumi}{1}
\tightlist
\item
  现在转到任意两个 \(v_i'\) 都相关的情形。令 \(A_i'\) 为 \(v_i' (1 \leqslant i \leqslant s)\) 的环；记 A 为赋值 \(v, A_i' \cap K = A\) 的环，对所有 \(i\) 都成立。令 \(B'\) 为由 \(A_1', \ldots, A_s'\) 生成的 L 的子环；记 \(B = B' \cap K\)；则 \(B \supset A\)。于是 B 是 K 上赋值 \(w\) 的环，而 \(B'\) 是赋值 \(w'\) 的环，该赋值非真且延拓 \(w\)（\(\S 7\)，第 2 节，命题 4）；域 \(\kappa(B')\) 是 \(\kappa(B)\) 的次数为 \(f(w'/w)\) 的扩张。考察典范像 \(\bar{A}_i', \bar{A}\)，它们是 \(A_i'\) 和 A 在 \(\kappa(B')\) 中的像；则 \(\bar{A}\) 是赋值 \(\bar{v}\) 在 \(\kappa(B)\) 上的环，而 \(\bar{A}_i'\) 是 \(\bar{v}_i'\) 在 \(\kappa(B')\) 上延拓 \(\bar{v}\) 的环。由于 \(A_i'\) 生成 \(B'\)，\(\bar{A}_i'\) 生成 \(\kappa(B')\)，所以 \(\bar{v}_i'\) 不全相关（\(\S 7\)，第 2 节，命题 4）。由证明的第一部分，
\end{enumerate}

\[
\sum_ {i = 1} ^ {s} e \left(\bar {v} _ {i} ^ {\prime} / \bar {v}\right) f \left(\bar {v} _ {i} ^ {\prime} / \bar {v}\right) \leqslant \left[ \kappa \left(\mathbf {B} ^ {\prime}\right): \kappa (\mathbf {B}) \right] = f \left(w ^ {\prime} / w\right)
\]

从而

\[
\sum_ {i = 1} ^ {s} e (w ^ {\prime} / w) e (\bar {v} _ {i} ^ {\prime} / \bar {v}) f (\bar {v} _ {i} ^ {\prime} / \bar {v}) \leqslant e (w ^ {\prime} / w) f (w ^ {\prime} / w) \leqslant n \quad \text {(no. 1, Lemma 1)}.
\]

因此，只要证明

\[
f (\bar {v} _ {i} ^ {\prime} / \bar {v}) = f (\bar {v} _ {i} ^ {\prime} / \bar {v}), \quad e (w ^ {\prime} / w) e (\bar {v} _ {i} ^ {\prime} / \bar {v}) = e (v _ {i} ^ {\prime} / v).\tag{8}
\]

为此，注意到 v 和 \(\bar{v}\)（分别地，\(v_{i}^{\prime}\) 和 \(\bar{v}_{i}^{\prime}\)）具有相同的剩余域（§ 4，第 1 节，命题 2 的推论）；这证明了第一个等式。对于第二个等式，根据 § 4，第 3 节的注，有如下交换图，其中各行是正合列，竖直箭头表示典范嵌入：

\[
\begin{array}{c} 0 \to \Gamma_ {\bar {v}} \to \Gamma_ {v} \to \Gamma_ {w} \to 0 \\ \downarrow \quad \downarrow \quad \downarrow \\ 0 \to \Gamma_ {\bar {v} _ {i} ^ {\prime}} \to \Gamma_ {v _ {i} ^ {\prime}} \to \Gamma_ {w ^ {\prime}} \to 0 \end{array}
\]

由此得到正合列

\[
0 \rightarrow \Gamma_ {\bar {v} _ {i} ^ {\prime}} / \Gamma_ {\bar {v}} \rightarrow \Gamma_ {v _ {i} ^ {\prime}} / \Gamma_ {v} \rightarrow \Gamma_ {w ^ {\prime}} / \Gamma_ {w} \rightarrow 0
\]

由第一章第 1 节第 4 节命题 2，得到一个正合列，这证明了第二个公式 (8)。为了完成定理 1 的证明，还需证明 (c)。若 \(v_i'\) 的环包含 \(v_j'\) 的环，则 \(\Gamma_{v_i'}\) 可识别为商群 \(\Gamma_{v_j'} / H\)，其中 \(H\) 是孤立子群（§ 4，第 3 节）。由于复合的典范映射

\[
\Gamma_ {v} \rightarrow \Gamma_ {v _ {j} ^ {\prime}} \rightarrow \Gamma_ {v _ {j} ^ {\prime}} / H = \Gamma_ {v _ {i} ^ {\prime}}
\]

是单射，\(H \cap \Gamma_{v} = \{0\}\)，从而 \(H = \{0\}\)（引理 3，第 1 节）。于是 \(v_{i}'\) 和 \(v_{j}'\) 等价，故 i = j。

注。环 \(A_{i}^{\prime}\)（对应赋值 \(v_{i}^{\prime}\)，其中 \(i \in I\)）的交 C 是 A 在 L 中的整闭包（\(\S 1\)，第 3 节，定理 3 的推论 3）；此外由 (c) 以及 \(\S 7\)，第 1 节，命题 1 和 2 可知，C 是半局部环，其极大理想是交 \(m_{i} = C \cap m(A_{i}^{\prime})\)，并且 \(A_{i}^{\prime} = C_{m}\) 对所有 \(i \in I\) 都成立。

\subsection{初始分歧指数}

定义 4. 设 G 是一个全序交换群，H 是 G 的有限指数子群。G 的由严格正元素组成且包含 H 的所有元素 \textgreater0 的优势子集的数目，称为 H 在 G 中的初始指数，记作 \(\varepsilon(\mathbf{G}, \mathbf{H})\)。

根据下述命题，这个初始指数是自然数：

命题 3. 在定义 4 的假设下，若 G 的严格正元素集没有最小元素，则对所有 H 都有 \(\varepsilon(G, H) = 1\)。如果 G 存在最小的 \(>0\) 元素，并记 G' 为由它生成的子群，则

\[
\varepsilon (\mathrm{G}, \mathrm{H}) = (\mathrm{G} ^ {\prime}: (\mathrm{G} ^ {\prime} \cap \mathrm{H})).
\]

在第一种情形中，令 x 为 G 中的一个元素 \textgreater0。满足 \(y \in G\) 且 0 \textless{} y \textless{} x 的集合是无限的，因而其中存在两个不同但模 H 同余的元素；它们的差是 H 中满足 0 \textless{} z \textless{} x 的元素。因此，每个包含 H 的全部严格正元素的优势子集都包含 x，进而包含 G 的所有元素 \textgreater0。

在第二种情形中，令 \(x\) 为最小的 \(>0\) 元素（在 \(G\) 中），并令 \(n\) 为最小的 \(>0\) 整数，使 \(nx \in H\)。显然 \(n = (G': (G' \cap H))\)。另一方面，记 \(M(y)\) 为满足 \(z \in G\) 且 \(y \leqslant z\)（\(y \in G\)）的 z 的集合，立即可见，定义 4 的优势集恰为 \(M(x), M(2x), \ldots, M(nx)\)。

推论。初始指数 \(\varepsilon(\mathbf{G},\mathbf{H})\) 整除指数 (G: H)，当 G 同构于 Z 时它等于该指数。

特别地，\(\varepsilon(G, H) \leqslant (G: H)\)。

定义 5. 设 K 为域，L 为 K 的有限扩张，w 为 L 上的赋值，v 为其在 K 上的限制，\(\Gamma_{w}\) 和 \(\Gamma_{v}\) 为它们的阶群。\(\Gamma_{v}\) 在 \(\Gamma_{w}\) 中的初始指数称为 w 关于 v（或关于 K）的初始分歧指数，记作 \(\varepsilon(w/v)\)。

由上述推论，\(\varepsilon(w/v)\) 整除 \(e(w/v)\)，在离散赋值的情形下二者相等。

命题 4. 在定义 5 的假设下，设 A 和 m（分别地，A' 和 m'）为赋值 v（分别地，w）的环和理想。则

\[
[ \mathrm{A} ^ {\prime} / \mathrm{mA} ^ {\prime}: \mathrm{A} / \mathrm{m} ] = \varepsilon (w / v) f (w / v).
\]

包含 mA' 且不同于 A' 的 A' 的理想，与由 \(\Gamma_{w}\) 的元素 \textgreater0 组成并包含 \textgreater0 的 \(\Gamma_{v}\) 的元素的优势子集相对应（\(\S3\)，第 5 节，命题 7 的推论）。因此它们的数目等于 \(\varepsilon(w/v)\)；由于它们关于包含关系构成全序集，这个数目等于商环 A'/mA' 的长度。现在，A' 上长度为 1 的模是 A'/m' 上的一维向量空间，因而是 A 上长度为 \(f(w/v)\) 的模；所以，既然 A'/mA' 作为 A'-模的长度为 \(\varepsilon(w/v)\)，它作为 A-模，也就是作为 A/m-模的长度为 \(\varepsilon(w/v)f(w/v)\)。

\subsection{\texorpdfstring{关系式 \(\sum_{i} e_{i} f_{i} = n\)}{关系式 Σ ei fi = n}}

命题 5. 设 K 为域，v 为 K 上的赋值，A 为其环，m 为其理想，L 为 K 的次数为 n 的有限扩张，B 为 A 在 L 中的整闭包，且 \((v_{i}^{\prime})_{1 \leq i \leq s}\) 是 v 到 L 的完备延拓系。则

\[
[ \mathrm{B} / \mathrm{mB}: \mathrm{A} / \mathrm{m} ] = \sum_ {i = 1} ^ {s} \varepsilon (v _ {i} ^ {\prime} / v) f (v _ {i} ^ {\prime} / v).
\]

令 \(A_{i}\) 为 \(v_{i}^{\prime}\) 的环；则 \(A_{i} = B_{m_{i}}\)，其中 \(m_{i}\) 遍历 B 的极大理想族（第 3 节，注）。令 \(q_{i}\) 为关于 \(m_{i}\) 的 mB 的饱和（第二章，第 2 节，第 4 节）。根据第五章第 2 节第 1 节命题 1 的推论 3，典范同态 \(B/mB \rightarrow \prod_{i=1}^{s} B/q_i\) 是同构，并且 \(m_i\) 是 B 中包含 \(q_i\) 的唯一极大理想。因此，\(B/q_i\) 典范同构于 \((\mathrm{B}/q_i)_{m_i}\)（第二章，第 3 节，第 3 节，命题 8），也就是

\[
\mathrm{B} _ {\mathrm{m} _ {i}} / \mathrm{mB} _ {\mathrm{m} _ {i}} = \mathrm{A} _ {i} / \mathrm{mA} _ {i}.
\]

因此存在典范同构 \(B/mB \rightarrow \prod_{i=1}^{s} A_{i}/mA_{i}\)，由第 4 节命题 4 即得结果。

推论。在相同假设和记号下，

\[
[ \mathrm{B} / \mathrm{mB}: \mathrm{A} / \mathrm{m} ] = \sum_ {i = 1} ^ {s} \varepsilon (v _ {i} ^ {\prime} / v) f (v _ {i} ^ {\prime} / v) \leqslant \sum_ {i = 1} ^ {s} e (v _ {i} ^ {\prime} / v) f (v _ {i} ^ {\prime} / v) \leqslant n.
\]

我们知道 \(\varepsilon(v_i'/v) \leqslant e(v_i'/v)\)（第 4 节，命题 3 的推论），并且 \(\sum_{i=1}^{s} e(v_i'/v)f(v_i'/v) \leqslant n\)（第 3 节，定理 1）。

定理 2. 在命题 5 的假设和记号下，下列条件等价：

\begin{enumerate}
\def\labelenumi{(\alph{enumi})}
\item
  B 是有限生成的 A-模；
\item
  B 是自由 A-模；
\item
  \([\mathbf{B} / \mathfrak{m}\mathbf{B}\colon \mathbf{A} / \mathfrak{m}] = n;\)
\end{enumerate}

\[
\text {(d)} \sum_ {i = 1} ^ {n} e (v _ {i} ^ {\prime} / v) f (v _ {i} ^ {\prime} / v) = n \text { and } \varepsilon (v _ {i} ^ {\prime} / v) = e (v _ {i} ^ {\prime} / v) \text { for   all } i.
\]

条件 (a) 与 (b) 的等价性由第 3 节第 6 节引理 1 得到。显然 (b) 蕴含 (c)（代数，第二章，第 1 节，第 5 节，公式 (19)）。(c) 与 (d) 的等价性由命题 5 的推论得到。还需证明 (c) 蕴含 (b)。

一般地，若 M 是 A-模，我们把 M/mM 关于 A/m 的向量空间记为 V(M)。假设 (c) 意味着 \(\dim(\mathrm{V}(\mathrm{B})) = n\)。令 \(x_{1}, \ldots, x_{n}\) 为 B 中的元素，它们在 V(B) 中的典范像构成 V(B) 的一个基，并令 \(L \subset B\) 为它们生成的 A-子模。由于 L 是无挠且有限生成的，它是自由的（§ 3，第 6 节，引理 1）。我们将看到 B = L。令 \(y \in B\)；记 \(M = L + Ay\)；这也是自由 A-模。典范嵌入 \(L \to M \to B\) 给出典范同态 \(\mathrm{V}(\mathrm{L}) \to \mathrm{V}(\mathrm{M}) \to \mathrm{V}(\mathrm{B})\)。由于 L 和 M 的秩均 \(\leqslant n\)，V(L) 和 V(M) 的维数也均不超过 n。现在由假设，\(\mathrm{V}(\mathrm{L}) \to \mathrm{V}(\mathrm{B})\) 是满射，而 V(B) 是 n 维的，故 V(L) 和 V(M) 都是 n 维的，且 \(\mathrm{V}(\mathrm{L}) \to \mathrm{V}(\mathrm{M})\) 是满射。由于 M 是有限生成的，\(L \to M\) 是满射（第二章，第 3 节，第 2 节，命题 4 的推论 1），从而 \(L = M, y \in L\)，且 B = L。因此 B 是自由的。

注 (1)。若 \(v\) 离散，\(\varepsilon(v_i'/v) = e(v_i'/v)\)（第 4 节），条件 (d) 化为 \(\sum_{i=1}^{s} e(v_i'/v)f(v_i'/v) = n\)。

推论 1. 在相同假设和记号下，再设 v 离散且 L 可分。则

\[
\sum_ {i = 1} ^ {s} e (v _ {i} ^ {\prime} / v) f (v _ {i} ^ {\prime} / v) = n.
\]

此时 A 的整闭包 B 是秩 n 的自由 A-模，因为 A 是主理想整环（第五章，第 1 节，第 6 节，命题 18 的推论 2）。

推论 2. 设 K 为域，v 是 K 上使 K 完备的离散赋值，L 为 K 的次数为 n 的有限扩张。~则 v 到 L 存在唯一（在等价意义下）的延拓 v'，v' 的环 A' 是 v 的环 A 上的有限生成自由模，且 \(e(v'/v)f(v'/v)=n\)。

\(v\) 到 \(L\) 的所有延拓都是相关的（第 2 节，命题 2 的推论）；由于它们是离散的（第 1 节，命题 1 的推论 3），所以它们等价（§ 4，第 5 节，命题 6 (c)）。这说明 \(v'\) 的唯一性。\(A\) 在 \(L\) 中的整闭包于是为 \(A'\)（§ 1，第 3 节，定理 3 的推论 3）。由于 \(v\) 离散，拓扑在 \(A\) 上由 \(K\) 上的拓扑诱导，是 m-进拓扑（其中 \(m = m(A)\)）；环 \(A\) 是完备的，因为它在 \(K\) 中闭。由于 \(A'/mA'\) 是有限维的向量 \((A/m)\)-空间（第 4 节，命题 4），\(A'\) 是有限生成的 A-模（第三章，第 2 节，第 9 节，命题 12 的推论 3）。因此由定理 2，A' 是自由的且 \(e(v'/v)f(v'/v) = n\)。

推论 3. 假设 \(v\) 的高为 1，并且定理 2 的等价条件成立；若 \(\hat{\mathbf{L}}_i\) 是 \(\mathbf{L}\) 关于 \(v_i'\) 的完备化，则次数 \(n_i = [\hat{\mathbf{L}}_i: \hat{\mathbf{K}}]\) 等于 \(e(v_i'|v)f(v_i'|v)\)，对所有 \(i\) 都成立，且典范同态

\[
\phi \colon \hat {\mathbf {K}} \otimes_ {\mathbf {K}} \mathbf {L} \rightarrow \prod_ {i = 1} ^ {s} \hat {\mathbf {L}} _ {i}
\]

（第 2 节，命题 2）。是双射的。对于所有 \(x \in \mathbf{L}\)，特征多项式 \(\mathrm{Pc}_{\mathrm{L} / \mathbb{K}}(x; \mathbf{X})\) 等于特征多项式 \(\mathrm{Pc}_{\hat{\mathrm{L}}_i / \hat{\mathbf{K}}}(x; \mathbf{X})\) 的乘积（\(1 \leqslant i \leqslant s\)）；特别地，

\[
\left\{ \begin{array} { l } \mathrm{Tr} _ { \mathbf { L } / \mathbb { K } } ( x ) = \sum _ { i = 1 } ^ { s } \mathrm{Tr} _ { \hat { \mathbf { L } } _ { i } / \hat { \mathbf { K } } } ( x ) \\ \mathrm{N} _ { \mathbf { L } / \mathbb { K } } ( x ) = \prod _ { i = 1 } ^ { s } \mathrm{N} _ { \hat { \mathbf { L } } _ { i } / \hat { \mathbf { K } } } ( x ) \\ v ( \mathrm{N} _ { \mathbf { L } / \mathbb { K } } ( x ) ) = \sum _ { i = 1 } ^ { s } n _ { i } v _ { i } ^ { \prime } ( x ) . \\ \end{array} \right.\tag{9}
\]

（(9) 中最后一个关系是有意义的，因为显然可以假设 \(v_{i}^{\prime}\)（由第 1 节命题 1 的推论 2 知其高为 1）像 \(v\) 一样将其值取在 \(\mathbf{R}\) 的一个子群中。）

由于 \(v_{i}^{\prime}\) 两两不等价且它们的高为 1，它们独立，因此第 2 节命题 2 表明 \(e(v_{i}^{\prime}/v)f(v_{i}^{\prime}/v) \leqslant n_{i}\) 对所有 \(i\) 成立，且 \(\sum_{i=1}^{s} n_i \leqslant n\)。所以第一个断言由这些不等式以及关系 \(\sum_{i=1}^{s} e(v_i'/v) f(v_i'/v) = n\) 得到。在同构 \(\phi\) 下，自同态 \(z \mapsto z(1 \otimes x)\) 作用于 \(\hat{\mathbf{K}} \otimes_{\mathbb{K}} \mathbf{L}\)（其中 \(x \in \mathbf{L}\)）时，变为 \(\prod_{i=1}^{s}\hat{L}_i\) 上的自同态；它保持每个因子不变，并在每个因子上化为乘以 \(x\)（因为 L 典范嵌入其完备化 \(\hat{L}_i\)）；由此得到关于 \(x\) 的特征多项式以及 (9) 的前两个公式的断言。最后，令 E 为 \(\hat{K}\) 的一个有限拟 Galois 扩张，包含 \(\hat{L}_i\)；由于 \(\hat{K}\) 完备且 \(\hat{v}\) 的高为 1，存在唯一一个作用于 E 的赋值 \(w\)（在等价意义下），它延拓 \(\hat{v}\)（第 2 节，命题 2 的推论 1）；于是对 E 的每个 \(\hat{K}\)-自同构 \(\sigma\)，都有 \(w(\sigma(x)) = v_i'(x)\)。因此

\[
\hat {v} \left(\mathrm{N} _ {\hat {\mathrm{L}} _ {i} / \hat {\mathrm{K}}} (x)\right) = n _ {i} v _ {i} ^ {\prime} (x)
\]

（代数，第 VIII 章，第 12 节，第 2 节，公式 (15)），这证明了 (9) 中的公式。

推论 4. 在推论 3 的假设下，若 L 是 K 的可分扩张，则每个 \(\hat{\mathbf{L}}_i\) 都是 \(\hat{\mathbf{K}}\) 的可分扩张。若进一步 L 是 K 的 Galois 扩张，Galois 群为 \(\mathcal{G}\)，且 \(\mathcal{G}_i\) 表示 B 中 \(v_i'\) 的理想的分解群（第五章，第 2 节，第 2 节，定义 2），则 \(\hat{\mathbf{L}}_i\) 是 \(\hat{\mathbf{K}}\) 的 Galois 扩张，其 Galois 群同构于 \(\mathcal{G}_i\)。

显然 \(\hat{\mathbf{L}}_i = \hat{\mathbf{K}}(\mathbf{L})\)；因此，若 \(\mathbf{L}\) 是 \(\mathbf{K}\) 的可分扩张，则 \(\hat{\mathbf{L}}_i\) 是 \(\hat{\mathbf{K}}\) 的可分扩张（代数，第五章，第 7 节，第 6 节，命题 10）。现在设 \(\mathbf{L}\) 是 Galois 扩张。每个 \(\sigma \in \mathcal{G}_i\) 都是 \(\mathbf{L}\) 上关于由 \(v_i'\) 定义的拓扑的连续自同构；由于 \(v_i'\) 的任意两个理想关于包含关系都不可比较（§ 7，第 2 节，定理 1 的推论 1），由此必有 \(v_i' = v_i' \circ \sigma\)，这由 \(\mathcal{G}_i\) 的定义得到；因此 \(\sigma\) 可通过连续性延拓为 \(\hat{\mathbf{K}}\)-自同构 \(\hat{\sigma}\)，作用于 \(\hat{\mathbf{L}}_i\)。这证明了 \(\hat{\mathbf{K}}\)-自同构在 \(\hat{\mathbf{L}}_i\) 上的数目至少等于 \(\operatorname{Card}(\mathcal{G}_i)\)。但由于赋值 \(v_i'\) 在 \(\mathcal{G}\) 下两两共轭（第五章，第 2 节，第 3 节，命题 6），有 \(s = (\mathcal{G}: \mathcal{G}_i)\)，从而

\[
\operatorname{Card} \left(\mathcal {G} _ {i}\right) = n / s \leqslant n.
\]

另一方面，根据推论 3，\(n = sn_{i}\)；这证明了 \(\hat{\mathbf{L}}_{i}\) 是 \(\hat{\mathbf{K}}\) 的 Galois 扩张，并且 \(\sigma \in \mathcal{G}_{i}\) 的自同构的连续延拓恰好构成 \(\hat{\mathbf{K}}\)-自同构在 \(\hat{\mathbf{L}}_{i}\) 上的全体。

注 (2)。上述部分结果可推广到 K 为不必交换的域上的赋值情形（参见~§ 3，第 1 节）。设 L 是 K 的扩域，且 \(v'\) 是 L 上的赋值，\(v\) 是其在 K 上的限制，A' 与 A 分别是赋值 \(v'\) 与 v 的环；则可如第 1 节定义分歧指数 \(e(v'/v)\)；另一方面，\(\kappa(\mathbf{A})\) 可视为 \(\kappa(\mathbf{A}')\) 的子域，并将 \(v'\) 关于 v 的（左）剩余秩定义为数 \(f(v'/v)\)，它等于左向量 \(\kappa(\mathbf{A})\)-空间 \(\kappa(\mathbf{A}')\) 的维数；若此维数有限，则如此定义，否则取 \(+\infty\)。于是，若 L 是关于 K 的有限维左向量空间，第 1 节引理 2 及其证明保持不变。此外，若 K 关于 v 完备，则第 5 节定理 2 的推论 2 的断言（\(v'\) 的存在性除外）也成立（n 表示 L 作为关于 K 的左向量空间的维数），证明如下：

首先，\(v'\) 在 \(L\) 上定义的拓扑是豪斯多夫的，并且与其左 \(K\)-向量空间结构相容，因此 \(v\) 到 \(L\) 的两个延拓在 \(L\) 上给出相同的拓扑（§ 5，第 2 节，命题 4），这证明这些延拓至多相差一个等价关系（§ 6，第 2 节）。下面证明：若 \(m = m(A)\)，则 \(A'/mA'\) 是左向量 \((A/m)\)-空间，其维数为 \(e(v'/v)f(v'/v)\)。记 \(e = e(v'/v)\)，可设 \(v(K^*) = Z\) 且 \(v'(L^*) = e^{-1}Z\)；令 \(u'\) 是 \(L\) 中满足 \(v'(u') = e^{-1}\) 的元素，令 \(u\) 是 \(K\) 中满足 \(v(u) = 1\) 的元素；于是 \(u = zu'^e\)，其中 \(z \in L\) 且 \(v'(z) = 0\)。由于 \(m\) 由 \(u\) 生成（作为 \(A\) 的左理想或右理想），\(mA' = u'^eA' = A'u'^e\)，只需证明：对于 \(0 \leqslant k \leqslant e - 1\)，\(A'u'^k/A'u'^{k+1}\) 是左向量 \((A/m)\)-空间，维数为 \(f(v'/v)\)。但 \(t \mapsto tu'^k\) 是从左 A-模 \(A'\) 到左 A-模 \(A'u'^k\) 的同构，将 \(A'u'\) 映到 \(A'u'^{k+1}\)，因而取商给出一个 \((A/m)\)-同构，将 \(A'/A'u'\) 映到 \(A'u'^k/A'u'^{k+1}\)；由 \(f(v'/v)\) 的定义以及 \(u'\) 生成 \(A'\) 的极大理想，便得到所述断言。证明与 \(K\)、\(L\) 交换时相同（有限生成无挠 A-模为自由模这一事实按 § 3，第 6 节，引理 1 的方式证明）。

\subsection{代数扩张中的赋值环}

命题 6. 设 K 为域，v 是 K 上的赋值，A 是其环，L 是 K 的代数扩张，A' 是 A 在 L 中的整闭包。设 V 为 L 上延拓 v 的赋值环的集合，M' 为 A' 的极大理想集合。则映射 V \(\mapsto\) m(V) \(\cap\) A' 是从 V 到 M' 的双射，而 \(m' \mapsto A'_{m'}\) 是逆双射。

A' 中的每个极大理想 \(m'\)（属于 \(A'\)）都满足 \(m' \cap A\) 是 A 的极大理想（第五章，第 2 节，第 1 节，命题 1），且 \(A_{m'}'\) 被 L 的某个赋值环 V 支配（因而它是 L 上延拓 v 的某个赋值的环）（§1，第 2 节，定理 2 的推论）。域 L 是由有限于 K 的 L 的子扩域 \(K_{\alpha}\) 组成的一个有向族的并；为了证明 \(V = A_{m'}'\)，只需证明 \(V \cap K_{\alpha} = A_{m'}' \cap K_{\alpha}\) 对所有 \(\alpha\) 都成立。现在，若记 \(A_{\alpha}' = A' \cap K_{\alpha}\)，则 \(A_{\alpha}'\) 是 A 在 \(K_{\alpha}\) 中的整闭包，因而是 \(K_{\alpha}\) 上延拓 v 的各赋值的环的交，而这些环 \(V_{i\alpha}\) 的数目有限，并且是局部环 \((\mathrm{A}_{\alpha}')_{\mathrm{m}'_{i\alpha}}\) 的 \(A_{\alpha}' (1 \leqslant i \leqslant n)\)，其中 \(m'_{i\alpha}\) 是 \(A_{\alpha}' (no. 3, Remark)\) 的不同极大理想；但 \(m' \cap A_{\alpha}'\) 是其中一个 \(m'_{i\alpha}\)，所以 \(V \cap K_{\alpha}\) 等于相应的局部环 \((\mathrm{A}'_{\alpha})_{m'_{i\alpha}} \subset \mathrm{A}'_{m'}\)，这就完成了 \(V = A'_{m'}\) 的证明。反之，若 \(V \in V\)，则 \(A' \subset V\)（§3，第 3 节，命题 6），且若 \(m' = m(V) \cap A'\)，则 \(m' \cap A = m\)，所以 \(m'\) 是 \(A'\) 的极大理想（第五章，第 2 节，第 1 节，命题 1），上述论证表明 \(V = A'_{m'}\)。

命题 7. 设 K 为域，L 是 K 的拟 Galois 扩张，f 和 \(f'\) 是取值于同一域 F 的 L 的位。假设 f 和 \(f'\) 在 K 上的限制相同。则存在 L 的一个 K-自同构 s，使得 \(f' = f \circ s\)。

设 A 为 K 上 f 与 \(f'\) 的公共限制所对应的位的环。f 与 \(f'\) 的环包含 A 在 L 中的整闭包 \(A'\)（§ 1，第 3 节，定理 3 的推论 3），因而（第五章，第 2 节，第 3 节，命题 6 的推论 1）存在 L 的一个 K-自同构 s，使得 \(f'\) 与 \(f \circ s\) 在 \(A'\) 上的限制相等；若 \(m'\) 是这些限制的公共核，则 \(m' \cap A\) 是 A 的极大理想，因而 \(m'\) 是 \(A'\) 的极大理想，且位 \(f'\) 与 \(f \circ s\) 在环 \(A_{m'}\) 上重合；但由命题 6，L 中支配 \(A_{m'}\) 的唯一赋值环就是环 \(A_{m'}\) 本身，所以位 \(f'\) 与 \(f \circ s\) 的环相同。

推论 1. 设 K 为域，v 是 K 上的赋值，L 是 K 的拟 Galois 扩张，v' 与 v'' 是 v 到 L 的两个延拓。则存在 L 的一个 K-自同构 s，使得 v'' 等价于 v' 与 s 的复合。

令 \(f'\) 和 \(f''\) 为分别与 \(v'\) 和 \(v''\) 相关的 K 的位；必要时以等价的位替换它们，可以设它们都将值取在 \(v\) 的剩余域的代数闭包中（第 1 节，命题 1）。于是存在 L 的一个 K-自同构 \(s\)，使得 \(f'' = f' \circ s\)（命题 7）；由位与赋值之间的对应（§ 3，第 3 节），\(v''\) 等价于 \(v' \circ s\)。

推论 2. 设 K 为域，f 是 K 的一个位（分别地，v 是 K 上的一个赋值），L 是 K 的纯不可分扩张。则 f（分别地，v）到 L 的所有延拓都等价。

L 是拟 Galois 扩张，且其唯一的自同构是恒等自同构。因此，推论 2 由命题 7（分别地，推论 1）得到。

命题 8. 设 K 为域，v 是 K 上的赋值，L 是 K 的次数为 n 的有限拟 Galois 扩张，且 \((v_{i}^{\prime})_{1\leq i\leq g}\) 是 v 到 L 的一个完备延拓系。则 \(e(v_{i}^{\prime}|v)\) 和 \(f(v_{i}^{\prime}|v)\) 的值 e、f 与 i 无关。于是 \(efg \leqslant n\)。若 v 的环 A 在 L 中的整闭包是有限生成的 A-模，则 \(efg = n\)。

这立即由定理 1（第 3 节）和定理 2（第 5 节）得到。

\subsection{绝对值的延拓}

命题 9. 设 K 为域，L 为 K 的代数扩张，f 为 K 上的绝对值。则 f 可延拓为 L 上的绝对值。

首先假设存在一个赋值 \(v\)，它在 \(\mathbf{K}\) 上取实值，使得 \(f(x) = e^{-v(x)}\)。存在一个赋值 \(v'\)，它在 \(\mathbf{L}\) 上，其在 \(\mathbf{K}\) 上的限制与 \(v\) 等价（§ 3，第 3 节，命题 5）。于是 \(v'\) 的高为 0 或 1（第 1 节，命题 1 的推论 2），因而可设其取实值。映射 \(x \mapsto e^{-v'(x)}\) 在 \(\mathbf{K}\) 上的限制是与 \(f\) 等价的绝对值，因而具有 \(f^s\) 的形式，其中 \(s > 0\)（一般拓扑，第 IX 章，第 3 节，第 2 节，命题 5）。于是

\[
x \mapsto e ^ {- v ^ {\prime} (x) / s}
\]

是 \(\mathbf{L}\) 上延拓 \(f\) 的绝对值。

现在假设 \(f\) 不是超度量绝对值。于是 \(K\) 可识别为 \(\mathbf{C}\) 的一个子域，使得 \(f(x) = |x|^s\)，其中 \(0 \leqslant s \leqslant 1\)（§ 6，第 4 节，定理 2）。由于 \(\mathbf{C}\) 代数闭，\(L\) 可识别为 \(\mathbf{C}\) 的一个子域，且绝对值 \(x \mapsto |x|^s\) 延拓 \(f\)。

命题 10. 设 K 为域，f 为 K 上的绝对值，K 关于 f 完备且非离散，L 为 K 的代数扩张。则 f 可唯一延拓为 L 上的绝对值 \(f'\)，并且若 L 的次数为 n，则

\[
f ^ {\prime} (x) = \left(f (\mathrm{N} _ {\mathrm{L/K}} (x))\right) ^ {1 / n}
\]

对所有 \(x \in \mathbf{L}\) 成立。

\(f'\) 的存在性由命题 9 得到，而其唯一性（在 \(L\) 的每个有限子扩张上，因而在整个 \(L\) 上）由 § 6，第 4 节的引理 2 得到。令 \(f'\) 为 \(f\) 到 \(K\) 的代数闭包的唯一延拓，并

假设 L 的次数为 n。~我们知道 \(\mathrm{N}_{\mathrm{L}/\mathrm{K}}(x)=\prod_{i=1}^{n} x_i\)，其中每个 \(x_{i}\) 都是 x 在 K 上的共轭元（代数，第 VIII 章，第 12 节，第 2 节，命题 4）。由唯一性，\(f', f'(x_{i})=f'(x)\) 对所有 i 成立，从而得到所述公式。

命题 11. 设 K 为域，f 为 K 上的非超度量绝对值，\(\hat{K}\) 为 K 关于 f 的完备化，\(\hat{f}\) 为 f 到 \(\hat{K}\) 的连续延拓，L 为 K 的次数为 n 的有限扩张。

\begin{enumerate}
\def\labelenumi{(\alph{enumi})}
\item
  令 \(f'\) 为 \(\mathbf{L}\) 上延拓 \(f\) 的绝对值；令 \(\hat{\mathbf{L}}_{f'}\) 表示 \(\mathbf{L}\) 关于 \(f'\) 的完备化，并将 \(\hat{\mathbf{K}}\) 识别为 \(\mathbf{K}\) 在 \(\hat{\mathbf{L}}_{f'}\) 中的闭包；则 \([\hat{\mathbf{L}}_{f'}: \hat{\mathbf{K}}] \leqslant n\)。
\item
  L 上延拓 f 的绝对值的数目有限。若将它们记为 \(f_{1}^{\prime}, \ldots, f_{s}^{\prime}\)，并将 L 关于 \(f_{i}^{\prime}\) 的完备化记为 \(\hat{L}_{i}\)，则典范映射
\end{enumerate}

\(\hat{\mathbf{K}}\otimes_{\mathbf{K}}\mathbf{L}\rightarrow \prod_{i = 1}^{n}\hat{\mathbf{L}}_i\) 是同构，并且

\[
\sum_ {i = 1} ^ {s} \left[ \hat {\mathrm{L}} _ {i}: \hat {\mathrm{K}} \right] = n.\tag{10}
\]

证明与命题 2 中相应断言的证明相同（第 2 节）。所引用的

§ 7，第 2 节，定理 1；§ 5，第 2 节，命题 4 的推论

应替换为下列

§ 7，第 3 节，定理 2；拓扑向量空间，第一章，

第 2 节，第 3 节，定理 2 的推论 1。

注意，\(f\) 到 \(\mathbf{L}\) 的两个定义相同拓扑的延拓相等（一般拓扑，第 IX 章，第 3 节，第 2 节，命题 5）。最后，由于 \(f\) 不是超度量绝对值，\(\mathbf{K}\) 的特征为 0，因而 \(\hat{\mathbf{K}} \otimes_{\mathbf{K}} \mathbf{L}\) 的 Jacobson 根为零。

\textbf{注}\par

\begin{enumerate}
\def\labelenumi{(\arabic{enumi})}
\item
  命题 11 (b) 表明，\(\hat{\mathbf{K}}\) 与 \(\mathbf{L}\) 在 \(\mathbf{K}\) 上的每个复合扩张都同构于某个完备化 \(\hat{\mathbf{L}}_i\)，并且这些复合扩张两两不同构。
\item
  我们知道完备化 \(\hat{\mathbf{K}}\) 和 \(\hat{\mathbf{L}}_i\) 同构于 \(\mathbf{R}\) 或 \(\mathbf{C}\)（§ 6，第 4 节，定理 2）。如果 \(\hat{\mathbf{K}}\) 同构于 \(\mathbf{C}\)，那么 \(\hat{\mathbf{L}}_i\) 对所有 \(i\) 也如此，且 (10) 表明延拓 \(f_i'\) 的数目等于 \(n\)。如果 \(\hat{\mathbf{K}}\) 同构于 \(\mathbf{R}\)（例如 \(\mathbf{K} = \mathbf{Q}\)），令 \(r_1\)（分别地，\(r_2\)）表示指标 \(i\) 的数目，其中 \(\hat{\mathbf{L}}_i\) 同构于 \(\mathbf{R}\)（分别地 \(\mathbf{C}\)）；则 (10) 可写为：
\end{enumerate}

\[
r _ {1} + 2 r _ {2} = n.\tag{11}
\]

命题 12. 设 K 为域，f 为 K 上的绝对值，L 为 K 的拟 Galois 扩张，且 \(f'\) 和 \(f''\) 是 f 到 L 的两个延拓。则存在 L 的一个 K-自同构 s，使得 \(f'' = f' \circ s\)。

若 \(f\) 是超度量绝对值，命题 7 的推论 1（第 6 节）表明，存在一个 K-自同构 \(s\)，作用于 \(L\)，使得 \(f''\) 和 \(f' \circ s\) 是等价绝对值；于是存在实数 \(a > 0\)，使得 \(f''(x) = (f'(s(x)))^a\)，对所有 \(x \in L\) 成立。若 \(f\) 不是非真绝对值，取 \(x \in K^*\) 使得 \(f(x) \neq 1\)，即可知 \(a = 1\)。若 \(f\) 是非真绝对值，则 \(f'\) 和 \(f''\) 也如此（命题 1 的推论 2，第 1 节），并且 \(s\) 可取为恒等自同构。

若 \(f\) 不是超度量绝对值，则存在 \(\mathbf{Q}\)-同构 \(u', u''\)，它们将 \(\mathbf{L}\) 同构到 \(\mathbf{C}\) 的子域中，并存在实指数 \(a' > 0\)、\(a'' > 0\)，使得 \(f'(x) = |u'(x)|^{a'}\)，并且

\[
f ^ {\prime \prime} (x) = | u ^ {\prime \prime} (x) | ^ {a ^ {\prime \prime}}
\]

对所有 \(x \in \mathbf{L}\) 成立（§ 6，第 4 节，定理 2）。取 \(x = 2\)，可见 \(a' = a''\)。\(u'\) 和 \(u''\) 在 \(\mathbf{K}\) 上的限制通过连续性延拓为同构 \(u_1\) 和 \(u_2\)，从 \(\hat{\mathbf{K}}\) 到 \(\mathbf{R}\)（分别地，\(\mathbf{C}\)）。于是 \(u_2 \circ \bar{u}_1^{-1}\) 是有赋值域 \(\mathbf{R}\)（分别地，\(\mathbf{C}\)）的自同构，因而是恒等自同构（分别地，是恒等自同构或自同构 \(c: \zeta \to \bar{\zeta}\)）。必要时将 \(u'\) 替换为 \(c \circ u'\)，可设 \(u'\) 和 \(u''\) 在 \(\mathbf{K}\) 上的限制相同。借助这个公共限制，把 \(\mathbf{K}\) 识别为 \(\mathbf{C}\) 的子域；\(u'\) 和 \(u''\) 是从 \(\mathbf{L}\) 到 \(\mathbf{C}\) 的 K-同构。由于 \(\mathbf{L}\) 是 \(\mathbf{K}\) 的拟 Galois 扩张，存在一个 K-自同构 \(s\)，作用于 \(\mathbf{L}\)，使得 \(u'' = u' \circ s\)；由于 \(a' = a''\)，立即得到 \(f'' = f' \circ s\)。

注 (3)。若 \(\hat{\mathbf{K}}\) 同构于 \(\mathbf{R}\)，命题 12 表明，\(\hat{\mathbf{L}}_t\)（\(\mathbf{L}\) 的所有完备化，采用命题 11 的记号）彼此同构。因此，采用上面注 2 的记号，要么 \(r_1 = n\) 且 \(r_2 = 0\)，要么 \(r_1 = 0\) 且 \(2r_2 = n\)。

\section{应用：局部紧域}

\subsection{局部紧域上的模函数}

设 K 为局部紧域（不必交换）。回顾，函数 mod（或 mod \(_K\)）已在 K 上定义如下（积分，第 VII 章，第 1 节，第 10 节，定义 6）：mod \(_K(0) = 0\)，对于 K 中 \(x \neq 0\)，数 mod \(_K(x)\) 是 K 的加法群的自同构 \(y \mapsto xy\) 的模。

命题 1. 若 K 是局部紧域，则函数 mod \(_{K}\) 属于 \(\mathcal{V}(K)\)（\(\S\) 6，第 1 节）。此外：

\begin{enumerate}
\def\labelenumi{(\roman{enumi})}
\item
  若 \(s > 0\) 使得 \((\mathrm{mod}_{\mathbf{K}})^s = g\) 是绝对值，则 \(g\) 定义 \(\mathbf{K}\) 上的拓扑。
\item
  若 K 非离散且 \(mod_{K}\) 是超度量绝对值，则 K 上存在一个赋范离散赋值 v，其环是紧的且剩余域是含 q 个元素的有限域，从而 \(mod_{K} = q^{-v}\)。K 上的拓扑由 v 定义。
\end{enumerate}

这由 § 6，第 1 节，命题 1；§ 5，第 1 节，命题 2；以及积分，第 VII 章，第 1 节，第 10 节，命题 12 和 13 得到。

命题 2. 设 K、K' 是两个（不必交换的）局部紧域，K 是 K' 的拓扑子域且 K 非离散。则：

\begin{enumerate}
\def\labelenumi{(\roman{enumi})}
\item
  \(\mathbf{K}'\) 是 \(\mathbf{K}\) 上的有限维左（分别地，右）向量空间。
\item
  如果 \(\mathbf{K}\) 包含于 \(\mathbf{K}'\) 的中心中，则对所有 \(x \in \mathbf{K}'\)，
\end{enumerate}

\[
\operatorname{mod} _ {\mathbf {K} ^ {\prime}} (x) = \operatorname{mod} _ {\mathbf {K}} (\mathrm{N} _ {\mathbf {K} ^ {\prime} / \mathbf {K}} (x)).\tag{1}
\]

由于 K 是非离散的完备赋值域，断言 (i) 由《拓扑向量空间》第一章，第 2 节，第 4 节，定理 3 得到；断言 (ii) 正是《积分》第七章，第 1 节，第 11 节，命题 17。

推论 1. 每个中心非离散的局部紧域都是其中心上的有限秩扩张。

局部紧域 K 的中心 Z 在 K 中是闭的，因而是局部紧的。

推论 2. 设 K' 是局部紧域，K 是 K' 的闭子域（两者均不必交换）。若 K' 是 K 上的 n 维左（分别地，右）向量空间，则

\[
\operatorname{mod} _ {\mathrm{K} ^ {\prime}} (x) = (\operatorname{mod} _ {\mathrm{K}} (x)) ^ {n} \quad \text {for all} \quad x \in \mathrm{K}.\tag{2}
\]

一般已知，在 K 上有限维的（左或右）向量空间中，比例为 \(x \in K\) 的伸缩的模等于 \((\text{mod}_{\mathbf{K}}(x))^{n}\)；只需将此应用于 \(K'\) 即可。

\subsection{代表元的存在}

命题 3. 设 K 为一个（不必交换的）非离散局部紧域，其拓扑由离散赋值 v 定义；令 A 为 v 的环，m 为 v 的理想，并记 \(\operatorname{Card}(\mathbf{A}/\mathfrak{m}) = q = p^{f}\)（p 为素数）。则存在 A/m 在 A 中的一个代表系 S 以及 v 的一个一致化元 u，使得 \(0 \in S\)，\(S^{*} = S \cap K^{*}\) 是 \(K^{*}\) 的循环子群，且 \(u^{-1} Su = S\)。此外，A 的每个元素都可唯一写成 \(\sum_{i=0}^{\infty} s_{i} u^{i}\) 的形式，其中 \(s_{i} \in S\)。

我们将使用下述引理：

引理 1. 设 \(x, y\) 是 \(A\) 中两个可交换的元素，使得 \(x - y \in \mathfrak{m}^j\)（\(j \geqslant 1\)）；则 \(x^{p^n} - y^{p^n} \in \mathfrak{m}^{j + n}\) 对每个整数 \(n \geqslant 0\) 都成立。

对 \(n\) 归纳即可将其归结为证明 \(n = 1\) 的情形。于是

\[
x ^ {p} - y ^ {p} = (x - y) \left(x ^ {p - 1} + x ^ {p - 2} y + \dots + y ^ {p - 1}\right);
\]

第二因子是 \(p\) 个彼此模 \(m\) 同余的项之和；由于 \(A / m\) 的特征为 \(p, p\)。\(1 \in m\) 在 \(A\) 中，因而

\[
x ^ {p - 1} + x ^ {p - 2} y + \dots + y ^ {p - 1} \in \mathfrak {m};
\]

由此 \(x^{p} - y^{p}\in \mathfrak{m}^{j + 1}\)

我们知道乘法群 \((\mathbf{A} / \mathfrak{m})^*\) 是一个具有 \(q - 1\) 个元素的循环群（代数，第 V 章，第 11 节，第 1 节，定理 1）；令 \(x\) 为 A 中该群一个生成元的代表；则 \(x^{q} - x \in \mathfrak{m}\)，由引理 1 得 \(x^{q^{n + 1}} - x^{q^n} \in \mathfrak{m}^{1 + fn}\)，因为 \(x^{q}\) 与 \(x\) 可交换。由此证明 \((x^{q^n})_{n \geqslant 0}\) 是 A 中的 Cauchy 序列；由于 A 紧且完备，该序列在 A 中有极限 \(s\)，显然满足 \(s \equiv x\)（模 \(\mathfrak{m}\)）且 \(s^q = s\)。由于 \(s \neq 0\)，\(s^{q - 1} = 1\)；更确切地说，\(s\) 是 A 中的一个本原 \((q - 1)\) 次单位根。显然，由 0 和各个幂 \(s^j\)（\(0 \leqslant j \leqslant q - 2\)）组成的集合 S 是 A 模 \(\mathfrak{m}\) 的各类的代表系，并且在 A 的乘法下不变。

现在令 \(a\) 为 \(v\) 的一致化元，并考察内自同构 \(y \mapsto a^{-1}ya\)（定义在 \(K\) 上）；它把 \(A\) 映到自身，把 \(m\) 映到自身，因而取商后定义域 \(A/m\) 的自同构；已知（代数，第 V 章，第 11 节，第 4 节，命题 5）这样的自同构具有 \(z \mapsto z^{p^r}\) 的形式，其中 \(0 \leqslant r \leqslant f - 1\)。于是 \(a^{-1}s^j a \equiv s^{jp^r} \pmod{m}\) 对 \(0 \leqslant j \leqslant q - 2\) 成立；由于 \(a \in m\) 且 \(s \notin m\)，这意味着 \(s^{-j}as^{jp^r} \equiv a \pmod{m^2}\)。

记

\[
u = \sum_ {j = 0} ^ {q - 2} s ^ {- j} a s ^ {j p ^ {r}}.
\]

于是 \(u \equiv (q - 1)a \equiv -a \pmod{m^2}\)，因为 \(p.1 \in m\)；我们得出 \(u\) 也是 \(v\) 的一致化元；此外

\[
s ^ {- 1} u s ^ {p ^ {r}} = u\tag{3}
\]

由此得到 \(u^{-1} \mathrm{S}u = \mathrm{S}\)。

最后，对于所有 \(x \in \mathbf{A}\)，存在唯一序列 \((s_i)\)（\(i \in \mathbf{N}\)），使得 \(s_i \in \mathbf{S}\) 对所有 \(i\) 都成立，并且 \(x \equiv \sum_{i=0}^{n} s_i u^i \pmod{\mathfrak{m}^{n+1}}\) 对所有 \(n \geqslant 0\) 都成立：关于 \(n\) 的归纳立即给出这一点，因为每个元素 \(t\) 属于 \(\mathfrak{m}^{n+1}\)，都满足形如 \(t \equiv t'u^{n+1} \pmod{\mathfrak{m}^{n+2}}\) 的关系，其中 \(t'\) 是 \(S\) 中唯一确定的元素。于是 \(x = \sum_{i=0}^{\infty} s_i u^i\)，族 \((s_i)\) 满足该关系且 \(s_i \in S\) 对所有 \(i\) 都成立，并被唯一确定。

\subsection{局部紧域的结构}

域 Q 关于 Q 上的非真绝对值的完备化 R 和 \(Q_{p}\) 是局部紧的（p 为任意素数）。另一方面，对于素数 p 的每个幂 \(q = p^{f}\)，形式幂级数域 \(\mathbf{F}_{q}((\mathrm{T}))\)（在有限域 \(F_{q}\) 上）的拓扑由 §3 第 4 节例 3 中定义的赋值给出，因而是局部紧的：因为赋值环 \(F_{q}[[T]]\) 的极大理想由 T 生成；我们知道这个环关于 (T)-进拓扑完备（第三章，第 2 节，第 6 节，命题 6），又由于剩余域 \(F_{q}\) 是有限的，§ 5 第 1 节的命题 2 证明了我们的断言。反之：

定理 1. 设 K 为一个（不必交换的）非离散局部紧域。

\begin{enumerate}
\def\labelenumi{(\roman{enumi})}
\item
  若 \(\mathbf{K}\) 的特征为 0 且 \(\mathrm{mod}_{\mathbf{K}}\) 不是超度量绝对值，则 \(\mathbf{K}\) 同构于域 \(\mathbf{R}\)、\(\mathbf{C}\) 或 \(\mathbf{H}\) 之一。
\item
  若 \(\mathbf{K}\) 的特征为 0 且 \(\mathrm{mod}_{\mathbf{K}}\) 是超度量绝对值，则 \(\mathbf{K}\) 是 \(p\)-进域 \(\mathbf{Q}_p\) 上有限秩的代数。
\item
  若 \(\mathbf{K}\) 的特征为 \(p \neq 0\)，则它同构于一个中心为形式幂级数域 \(\mathbf{F}_q((\mathrm{T}))\) 的域（其中 \(q\) 是 \(p\) 的幂），并且在其中心上有限秩。
\end{enumerate}

\begin{enumerate}
\def\labelenumi{(\roman{enumi})}
\item
由 Ostrowski 定理（§ 6，第 4 节，定理 2），K 是一个同构于 R、C 或 H 的处处稠密子域的拓扑域；又由于 K 完备，它同构于 R、C 或 H。
\item
  令 A 为绝对值 \(mod_{K}\) 的环，m 为其极大理想。我们知道 A/m 是有限域（§ 5，第 1 节，命题 2），因而 \(mod_{K}\) 在 Q 上诱导的绝对值具有有限剩余域；这只有当它等价于 \(p\)-进绝对值时才可能（§ 6，第 3 节，命题 4）。\(\mathbf{Q}\) 在 \(\mathbf{K}\) 中的闭包于是同构于 \(\mathbf{Q}_p\)，并包含于 \(\mathbf{K}\) 的中心中，因为后者在 \(\mathbf{K}\) 中闭；由第 1 节的命题 2 即得结论。
\item
  第二个断言由第 2 节命题 2 的推论和第一个断言得到。为证明第一个断言，注意 \(\mathrm{mod}_{\mathbf{K}}\) 必为超度量绝对值（§ 6，第 2 节，命题 3 的推论）；在第 2 节命题 3 的证明所用的记号中，中心 \(Z\) 是 \(K\) 中与 \(s\) 和 \(u\) 都交换的元素组成的；但是根据 (3)，
\end{enumerate}

\[
u ^ {- q} s u ^ {q} = s ^ {q p ^ {r}} = s,
\]

于是 \(u^q \in \mathbf{Z}\)，我们得出 \(\mathbf{Z}\) 非离散。由于 \(\mathbf{Z}\) 局部紧，我们可限于 \(\mathbf{K}\) 交换的情形。此时的 \(\mathbf{F}_p\)-代数 \(\mathbf{F}_p[s]\) 位于 \(\mathbf{K}\) 中，是有限域，因为 \(s^{q-1} = 1\)，且显然其中每个元素都满足 \(y^q = y\)，所以它与 \(S\) 相同并同构于 \(\mathbf{F}_q\)，因为 \(S \subset \mathbf{F}_p[s]\) 有 \(q\) 个元素。由于 \(S\) 中任意两个元素之和仍在 \(S\) 中，把每个形式幂级数 \(\sum_{i=0}^{\infty} s_i T^i \in \mathbf{F}_q[[T]]\) 映到元素 \(\sum_{i=0}^{\infty}s_{i}u^{i}\) 的映射，是从环 \(F_{q}[[T]]\) 到环 A 的双射同态，从而立即得到结论。

推论 1. 每个非离散局部紧域在其中心上有限秩。

推论 2. 每个局部紧域或者连通，或者全不连通；若它连通，则同构于 R、C 或 H。

若域 K 上的拓扑由超度量绝对值定义，则 K 关于此拓扑全不连通。

注。令 \(s\) 为一个整数 \(>0\)；子域 \(\mathbf{F}_{q}((\mathrm{T}^{s})) = \mathbf{L}\) 是 \(\mathbf{K} = \mathbf{F}_{q}((\mathrm{T}))\) 的子域，在 \(\mathbf{K}\) 中闭，且 \(e(\mathrm{K/L}) = s\)、\(f(\mathrm{K/L}) = 1\)。因此可见，存在闭子域 \(\mathbf{L}\)，它们属于 \(\mathbf{K}\)，非离散且使 \(e(\mathrm{K/L})\)（从而次数 {[}K: L{]}）任意大（这与特征为 0 的局部紧域的情形相反；在那里，每个这样的局部紧子域 \(\mathbf{L}\) 都属于域 \(\mathbf{K}\)，且必包含 \(\mathbf{R}\) 或 \(\mathbf{Q}_{p}\)，因而 {[}K: L{]} 有界）。

\section{赋值到超越扩张的延拓}

\subsection{单生成超越扩张的情形}

引理 1. 设 K 为域，v 为 K 上的赋值，Γ 为其阶群，Γ' 是一个全序的

群，包含 \(\Gamma\)，且 \(\xi\) 是 \(\Gamma'\) 的一个元素。存在唯一的赋值 \(w\)，定义在 \(\mathbf{K}(\mathbf{X})\) 上，使得对于 \(\mathbf{P} = \sum_{j} a_j \mathbf{X}^j (a_j \in \mathbf{K})\)，有 \(w(\mathbf{P}) = \inf_j(v(a_j) + j\xi)\)。

根据 § 3 第 2 节的命题 4，只需证明公式

\[
w \left(\sum_ {j} a _ {j} \mathbf {X} ^ {j}\right) = \inf _ {j} (v (a _ {j}) + j \xi)\tag{1}
\]

在环 K{[}X{]} 上定义一个赋值。由于

\[
v \left(a _ {j} + b _ {j}\right) + j \xi \geqslant \inf (v \left(a _ {j}\right), v \left(b _ {j}\right)) + j \xi = \inf (v \left(a _ {j}\right) + j \xi , v \left(b _ {j}\right) + j \xi),
\]

由此可得

\[
w (\mathbf {P} + \mathbf {Q}) \geqslant \inf (w (\mathbf {P}), w (\mathbf {Q}))\tag{2}
\]

对于 \(\mathbf{P},\mathbf{Q}\) 属于 \(\mathbf{K}[\mathbf{X}]\)，当 \(w(\mathbf{P})\neq w(\mathbf{Q})\) 时取等号。我们证明

\[
w (\mathrm{PQ}) = w (\mathrm{P}) + w (\mathrm{Q})\tag{3}
\]

对于 \(\mathbf{P} = \sum_{j} a_{j} \mathbf{X}^{j}\) 和 \(\mathbf{Q} = \sum_{j} b_{j} \mathbf{X}^{j}\)。令 \(i\)（分别地，\(k\)）为下述整数 \(j\) 中的最小者：它们使 \(v(a_{j}) + j\xi\)（分别地，\(v(b_{j}) + j\xi\)）取得最小值；令 \(\alpha\)（分别地，\(\beta\)）表示这个最小值。对于 \(j, j'\) 属于 \(\mathbf{N}\)，

\[
w (a _ {j} b _ {j ^ {\prime}} \mathbf {X} ^ {j + j ^ {\prime}}) = v (a _ {j}) + j \xi + v (b _ {j ^ {\prime}}) + j ^ {\prime} \xi \geqslant \alpha + \beta ,
\]

由此根据 (2) 得 \(w(\mathrm{PQ}) \geqslant \alpha + \beta\)。现在考虑项 \(c\mathbf{X}^{i + k}\)，它在 PQ 中的次数为 \(i + k\)；则 \(c = \sum_{n \in \mathbf{Z}} a_{i + n} b_{k - n}\)；根据 \(i\) 和 \(k\) 的选取，该元素

\[
w (a _ {i + n} b _ {k - n} \mathbf {X} ^ {i + k}) = v (a _ {i + n}) + (i + n) \xi + v (b _ {k - n}) + (k - n) \xi
\]

其最小值 \(\alpha + \beta\) 仅在 \(n = 0\) 时且只在此时取得；因此 \(w(c\mathbf{X}^{i + k}) = \alpha + \beta\)，从而根据 (1)，

\[
w (\mathrm{PQ}) = \alpha + \beta = w (\mathrm{P}) + w (\mathrm{Q}).
\]

命题 1. 设 K 为域，v 为 K 上的赋值，\(\Gamma\) 为其阶群，\(\Gamma'\) 为包含 \(\Gamma\) 的全序群，且 \(\xi\) 是 \(\Gamma'\) 的一个元素，并满足关系 \(n\xi\in\Gamma\)，\(n\in Z\) 蕴含 n=0。则存在唯一的赋值 w，取值于 \(\Gamma'\) 且延拓 v，并满足 \(w(\mathbf{X})=\xi\)。w 的剩余域等于 v 的剩余域，其阶群是 \(\Gamma+Z\xi\)（\(\Gamma'\) 的子群）。

我们先证明 \(w\) 的唯一性。设 \(\mathbf{P} = \sum_{j} a_j \mathbf{X}^j\) 是 \(\mathbf{K}[\mathbf{X}]\) 中的元素。则 \(w(a_j \mathbf{X}^j) = v(a_j) + j\xi\)，这说明单项式 \(a_j \mathbf{X}^j\) 在 \(a_j \neq 0\) 时关于 \(w\) 的值彼此不同。由此 \(w(\mathbf{P}) = \inf_j (v(a_j) + j\xi)\)，这同时说明 \(w\) 在 \(\mathbf{K}[\mathbf{X}]\) 上（从而也在 \(\mathbf{K}(\mathbf{X})\) 上）的唯一性，并说明 \(w\) 的阶群是 \(\Gamma + \mathbf{Z}\xi\)。还可看出，若 \(\mathbf{P} \neq 0\)，则可写成 \(\mathbf{P} = a\mathbf{X}^n (1 + u)\) 的形式，其中 \(a\in \mathbf{K}^*\)，\(n\in \mathbf{N}\)，\(u\in \mathbf{K}(\mathbf{X})\)，且 \(w(u) > 0\)；因此，对于任意 \(\mathbf{R}\neq 0\)，它作为 \(\mathbf{K}(\mathbf{X})\) 中的元素都可写成

\[
\mathrm{R} = b \mathrm{X} ^ {n} (1 + u ^ {\prime}),
\]

其中 \(b \in \mathbf{K}^*\)，\(n \in \mathbf{Z}\)，\(u' \in \mathbf{K}(\mathbf{X})\)，且 \(w(u') > 0\)；于是 \(w(\mathbf{R}) = v(b) + n\xi\)，所以 \(w(\mathbf{R}) = 0\) 当且仅当 \(v(b) = 0\) 且 \(n = 0\)；因此，当 \(w(\mathbf{R}) = 0\) 时，\(\mathbf{R}\) 与 \(b\) 模 \(w\) 的理想同余，这说明 \(w\) 的剩余域等于 \(v\) 的剩余域。

最后，w 的存在性由引理 1 得出。

命题 2. 设 K 为域，v 为 K 上的赋值，\(\Gamma\) 为其阶群，k 为其剩余域。存在 K(X) 上唯一的延拓 v 的赋值 w，使得 \(w(X) = 0\)，且 X 在剩余域 \(k'\)（w 的剩余域）中的像 t 超越于 k。w 的阶群等于 v 的阶群，其剩余域为 \(k(t)\)。

为证明 \(w\) 的唯一性，只需证明：若 \(\mathbf{P} = \sum_{j} a_j \mathbf{X}^j\) 是 \(\mathbf{K}[\mathbf{X}]\) 中的非零元素，则

\[
w (\mathbf {P}) = \inf _ {j} (v (a _ {j})).
\]

可以用 \(K^{*}\) 中的一个元素去除 P，并设对所有 j 都有 \(v(a_{j}) \geqslant 0\)，且其中一个 \(v(a_{j})\) 为 0。由于 \(w(\mathbf{X}) = 0\)，P 于是属于 w 的环；记 \(\bar{a}_{j}\) 为 \(a_{j}\) 在 k 中的典范像，则 P 在剩余域 \(k'\) 中的典范像为 \(\sum_{j} \bar{a}_{j} t^{j}\)；由于 t 超越于 k 且其中一个 \(\bar{a}_{j}\) 非零，该像非零，从而

\[
w (\mathbf {P}) = 0 = \inf _ {j} (v (a _ {j})).
\]

现在证明 \(w\) 的存在性。根据引理 1，公式 \(w(\mathbf{P}) = \inf_{j}(v(a_j))\)（对于 \(\mathbf{P} = \sum_{j} a_j \mathbf{X}^j\)）定义赋值 \(w\)，作用在 \(\mathbf{K}(\mathbf{X})\) 上，且显然 \(w\) 与 \(v\) 有相同的阶群。于是 \(w(\mathbf{X}) = 0\)。我们证明 \(t\)（\(\mathbf{X}\) 在剩余域 \(k'\)（\(w\) 的剩余域）中的典范像）超越于 \(k\)：若 \(\sum_{j} \bar{a}_j t^j = 0\)，其中 \(\bar{a}_j \in k\) 对所有 \(j\) 成立，则取 \(a_j\) 为 \(\bar{a}_j\) 在 \(v\) 的环中的代表，有 \(w\left(\sum_{j} a_j \mathbf{X}^j\right) > 0\)；从而 \(v(a_j) > 0\) 对所有 \(j\) 成立，于是 \(\bar{a}_j = 0\) 对所有 \(j\) 成立。最后证明 \(k' = k(t)\)：每个元素 \(\mathbf{R}\) 都属于 \(\mathbf{K}(\mathbf{X})\)，且都可写成 \(\mathbf{R} = c\left(\sum_{j} a_j \mathbf{X}^j\right) / \left(\sum_{j} b_j \mathbf{X}^j\right)\) 的形式，其中 \(c, a_j, b_j\) 属于 \(\mathbf{K}\)，\(v(a_j) \geqslant 0\) 且 \(v(b_j) \geqslant 0\) 对所有 \(j\) 成立，其中一个 \(v(a_j)\) 和一个 \(v(b_j)\) 为零；于是 \(w(\mathbf{R}) \geqslant 0\) 当且仅当 \(v(c) \geqslant 0\)；记 \(f\) 为从 \(w\) 的环到 \(k'\) 的典范同态，

\[
f (\mathrm{R}) = f (c) \left(\sum_ {j} f (a _ {j}) t ^ {j}\right) / \left(\sum_ {j} f (b _ {j}) t ^ {j}\right),
\]

这就证明了我们的断言。

注。不要以为我们刚遇到的 v 到 K(X) 的两类延拓是仅有的两类；还可能存在第三类延拓，其中 \(\Gamma'/\Gamma\) 是扭群，而 \(k'\) 是 k 的（不必有限的）代数扩张。这第三类不一定由引理 1 中所述的步骤得到（参见~§3，习题 1）。

\subsection{交换群的有理秩}

定义 1. 交换群 G 的有理秩是向量 Q-空间 \(G \otimes_{\mathbf{Z}} \mathbf{Q}\) 的维数。

这个维数也可定义为下述基数 r 的最小上界（有限或无限）：存在 G 的 r 个元素，它们在 Z 上线性无关（《代数》第二章 § 7 第 10 节命题 26）。当且仅当 G 是扭群时，G 的有理秩为零。对于加法群 \(R^{n}\) 的子群，有理秩的概念与《一般拓扑学》第七章 § 1 中定义的概念一致。

在本段其余部分，我们用 \(r(G)\) 表示交换群 \(G\) 的有理秩。若 \(G'\) 是 \(G\) 的子群，则（由于 \(Q\) 是平坦的 \(Z\)-模）有加法等式

\[
r (\mathbf {G}) = r \left(\mathbf {G} ^ {\prime}\right) + r \left(\mathbf {G} / \mathbf {G} ^ {\prime}\right).\tag{4}
\]

命题 3. 设 G 是全序交换群，H 是 G 的子群。若 \(h(\mathbf{G})\) 和 \(h(\mathbf{H})\) 分别表示 G 和 H 的高（§ 4，第 4 节），则不等式

\[
h (\mathbf {G}) \leqslant h (\mathbf {H}) + r (\mathbf {G} / \mathbf {H})\tag{5}
\]

成立。

事实上，设 \(G_{0} \subset G_{1} \subset \cdots \subset G_{n}\) 是 G 的严格递增孤立子群序列。必须证明不等式

\[
n \leqslant h (\mathrm{H}) + r (\mathrm{G} / \mathrm{H}).\tag{6}
\]

当 n = 0 时显然。设 \(n \geqslant 1\)，并对 n 作归纳。~应用归纳假设于群 \(G_{n-1}\) 及其子群 \(H \cap G_{n-1}\)，得到

\[
n - 1 \leqslant h (\mathrm{H} \cap \mathrm{G} _ {n - 1}) + r (\mathrm{G} _ {n - 1} / (\mathrm{H} \cap \mathrm{G} _ {n - 1})).\tag{7}
\]

然后我们区分两种情形：

\begin{enumerate}
\def\labelenumi{(\alph{enumi})}
\tightlist
\item
  \(\mathbf{H} \cap \mathbf{G}_{n-1} = \mathbf{H}\)，换言之，\(\mathbf{H} \subset \mathbf{G}_{n-1}\)。不等式 (7) 可写成
\end{enumerate}

\[
n \leqslant h (\mathrm{H}) + r (\mathrm{G} _ {n - 1} / \mathrm{H}) + 1.\tag{8}
\]

现在 \(\mathbf{G} / \mathbf{G}_{n - 1}\) 是一个不化为 0 的全序群；因而它不是扭群，且 \(r(\mathrm{G} / \mathrm{G}_{n - 1}) \geqslant 1\)。由 (4) 得 \(r(\mathrm{G} / \mathrm{H}) \geqslant r(\mathrm{G}_{n - 1} / \mathrm{H}) + 1\)。代入 (8)，我们必然得到所需不等式 (6)。

\begin{enumerate}
\def\labelenumi{(\alph{enumi})}
\setcounter{enumi}{1}
\tightlist
\item
  \(\mathbf{H} \cap \mathbf{G}_{n-1} \neq \mathbf{H}\)。由于 \(\mathbf{H} \cap \mathbf{G}_{n-1}\) 是 \(\mathbf{H}\) 的孤立子群，我们得出 \(h(\mathbf{H}) \geqslant h(\mathbf{H} \cap \mathbf{G}_{n-1}) + 1\)。另一方面，显然 \(r(\mathbf{G}/\mathbf{H}) \geqslant r(\mathbf{G}_{n-1}/(\mathbf{H} \cap \mathbf{G}_{n-1}))\)。代入 (7)，再次得到 (6)。
\end{enumerate}

推论。对于每个全序交换群 \(\mathbf{G}\)，有 \(h(\mathbf{G}) \leqslant r(\mathbf{G})\)。

在命题 3 中令 \(\mathbf{H} = \{0\}\)。

命题 4. 设 G 是全序交换群。设 G 有限生成，且 \(h(\mathbf{G}) = r(\mathbf{G})\)。则 G 同构于按字典序排列的 \(\mathbf{Z}^{r(\mathbf{G})}\)。

记 \(r = r(\mathbf{G}) = h(\mathbf{G})\)。若 \(r = 0\)，则 \(\mathbf{G} = \{0\}\)。若 \(r = 1\)，有限生成交换群的结构表明，存在一个同构 \(j\)，它把 \(\mathbf{G}\) 映到 \(\mathbf{Z}\) 上（《代数》第七章 § 4 第 6 节定理 3）。现在，\(\mathbf{Z}\) 只有两个与其群结构相容的全序，即通常序和反序。因此，\(j\) 或 \(-j\) 是把有序群 \(\mathbf{G}\) 同构到带通常序的 \(\mathbf{Z}\) 上的同构。

现在设 \(r \geqslant 2\)，并对 \(r\) 作归纳。设 \(H\) 是 \(G\) 的高为 \(r - 1\) 的孤立子群。则 \(r(H) + r(G/H) = r\)（公式 (4)），\(r(H) \geqslant h(H) = r - 1\)，且 \(r(G/H) \geqslant h(G/H) = 1\)（命题 3 的推论），从而 \(r(H) = r - 1\) 且 \(r(G/H) = 1\)。归纳假设表明 \(H\) 同构于按字典序排列的 \(Z^{r-1}\)，而 \(r = 1\) 的情形表明 \(G/H\) 同构于 \(Z\)。由于 \(Z\) 是自由 \(Z\)-模，\(H\) 是 \(G\) 的直因子（《代数》第二章 § 1 第 11 节命题 21）。于是下述引理表明，\(G\) 同构（非典范地）于字典积 \(H \times (G/H)\)，从而完成证明。

引理 2. 设 H 是全序交换群 G 的孤立子群。若 H 是 G 的直因子，则有序群 G 同构于按字典序排列的群 (G/H) × H。

令 \(j\) 为从 \((\mathrm{G} / \mathrm{H}) \times \mathrm{H}\) 到 \(\mathbf{G}\) 的同构，使得 \(j(0, x) = x\) 对所有 \(x \in \mathrm{H}\) 成立，并且 \(j(y, x)\) 属于陪集 \(y\) 的 \(\mathrm{H}\)。由于 \((\mathrm{G} / \mathrm{H}) \times \mathrm{H}\) 是全序的，只需证明 \(j\) 是递增的（《集合论》第三章 § 1 第 12 节命题 11）。令 \((y, x)\) 为满足 \(\geqslant 0\) 的元素，它属于按字典序排列的 \((\mathrm{G} / \mathrm{H}) \times \mathrm{H}\)。若 \(y > 0\)，\(\mathrm{H}\) 的包含 \(j(y, x)\) 的陪集是一个 \(> 0\) 的元素，从而 \(j(y, x) > 0\)，因为否则 \(y \leqslant 0\)（§ 4 第 2 节命题 3）。若 \(y = 0\) 且 \(x \geqslant 0\)，则 \(j(y, x) = x \geqslant 0\)。因此 \(j\) 确为递增的。

\subsection{任意超越扩张的情形}

在本号中，我们使用下述记号：K 是域，\(K'\) 是 K 的扩域，v 是 K 上的赋值，\(v'\) 是 v 到 \(K'\) 的延拓，\(\Gamma\) 和 k（分别地，\(\Gamma'\) 和 \(k'\)）分别是 \(v\)（分别地，\(v'\)）的阶群和剩余域。我们记：

\[
d \left(\mathrm{K} ^ {\prime} / \mathrm{K}\right) = \dim . \mathrm{al} _ {\mathrm{K}} \mathrm{K} ^ {\prime} = \text { transcendence   degree   of } \mathrm{K} ^ {\prime} \text { over } \mathrm{K};
\]

\[
s (v ^ {\prime} / v) = \dim . \mathrm{al} _ {k} k ^ {\prime} = \text { transcendence   degree   of } k ^ {\prime} \text { over } k;
\]

\[
r (v ^ {\prime} / v) = r (\Gamma^ {\prime} / \Gamma) = \text { rational   rank   of } \Gamma^ {\prime} / \Gamma ,
\]

若等号右边各项有限；否则约定 \(d(\mathbf{K}' / \mathbf{K}) = +\infty\)（分别地，\(s(v' / v) = +\infty\)、\(r(v' / v) = +\infty\)）。

定理 1. 设 \(x_1, \ldots, x_s\) 是 \(v'\) 的环中的元素，其典范像 \(\bar{x}_i\) 在 \(k'\) 中在 \(k\) 上代数独立；设 \(y_1, \ldots, y_r\) 是 \(K'\) 中的元素，使得 \(v'(y_j)\) 在 \(\Gamma'/\Gamma\) 中的典范像在 \(Z\) 上线性无关。则 \(r + s\) 个元素 \(x_1, \ldots, x_s, y_1, \ldots, y_r\) 属于 \(K'\)，并在 \(K\) 上代数独立；\(v'\) 限制于 \(K(x_1, \ldots, x_s, y_1, \ldots, y_r)\) 后的剩余域为 \(k(\bar{x}_1, \ldots, \bar{x}_s)\)，并且

\[
\Gamma + \mathbf {Z} v ^ {\prime} (y _ {1}) + \dots + \mathbf {Z} v ^ {\prime} (y _ {r})
\]

作为有序群。

当 \(r + s = 0\) 时断言显然。我们对 \(r + s\) 作归纳。若 \(r' \leqslant r\)、\(s' \leqslant s\) 且 \(r' + s' < r + s\)，则归纳假设表明：当 \(K\) 换成 \(K(x_1, \ldots, x_s, y_1, \ldots, y_{r'})\)，且族 \((x_1, \ldots, x_s), (y_1, \ldots, y_r)\) 换成 \((x_{s'+1}, \ldots, x_s), (y_{r'+1}, \ldots, y_r)\) 时，定理 1 的假设仍成立。因此问题归结为下述两种情形之一：

\begin{enumerate}
\def\labelenumi{(\alph{enumi})}
\item
  存在 \(x\) 属于 \(v'\) 的环，使得 \(\bar{x}\) 在 \(k\) 上超越；则须证明 \(x\) 在 \(\mathbf{K}\) 上超越，且 \(v'\) 限制于 \(\mathbf{K}(x)\) 后的剩余域为 \(k(\bar{x})\)，阶群为 \(\Gamma\)。
\item
  存在 \(y\) 属于 \(\mathbf{K}'\)，使关系 \(nv'(y) \in \Gamma\) 和 \(n \in \mathbf{Z}\) 蕴含 \(n = 0\)；则须证明 \(y\) 在 \(\mathbf{K}\) 上超越，且 \(v'\) 限制于 \(\mathbf{K}(y)\) 后的剩余域为 \(k\)，阶群为 \(\Gamma + \mathbf{Z}v'(y)\)。
\end{enumerate}

现在第 8 节第 1 节的命题 1 表明，\(x\)（分别地，\(y\)）不可能在 \(K\) 上代数。(a)（分别地，(b)）的其余断言由第 1 节的命题 2（分别地，命题 1）立即得到。

推论 1. 不等式

\[
s (v ^ {\prime} / v) + r (v ^ {\prime} / v) \leqslant d (\mathbf {K} ^ {\prime} / \mathbf {K})\tag{9}
\]

成立。

此外，若 \(K'\) 是 K 的有限生成扩张，且 (9) 中取等号，则 \(\Gamma'/\Gamma\) 是有限生成 Z-模，且 \(k'\) 是 k 的有限生成扩张。

令 \(r\) 和 \(s\) 为自然数，满足 \(r \leqslant r(v'/v)\) 且 \(s \leqslant s(v'/v)\)；我们证明 \(r + s \leqslant d(\mathbf{K}'/\mathbf{K})\)，这将证明 (9)。根据假设，存在 \(x_1, \ldots, x_s, y_1, \ldots, y_r\) 属于 \(\mathbf{K}'\) 的元素，它们满足定理 1 的假设。

所以它们在 K 上代数独立，这就表明不等式 \(r + s \leqslant d(\mathbf{K}'/\mathbf{K})\) 成立。

若 \(\mathbf{K}'\) 是 \(\mathbf{K}\) 的有限生成扩张，\(d(\mathbf{K}' / \mathbf{K})\) 有限，因而 \(s(v' / v)\) 和 \(r(v' / v)\) 也有限；我们分别记为 \(s\) 和 \(r\)。存在元素 \(x_{1},\ldots ,x_{s},\) \(y_{1},\ldots ,y_{r}\) 属于 \(\mathbf{K}'\)，它们满足定理 1 的假设。若 \(r + s = d(\mathbf{K}' / \mathbf{K})\)，这些元素构成 \(\mathbf{K}'\) 在 \(\mathbf{K}\) 上的超越基，因而 \(\mathbf{K}'\) 是 \(\mathbf{K}'' = \mathbf{K}(x_1,\dots,y_r)\) 的有限代数扩张。令 \(\Gamma''\) 和 \(k''\) 为 \(v'\) 限制于 \(\mathbf{K}''\) 所得赋值的阶群和剩余域。根据定理 1，\(\Gamma'' / \Gamma\) 是有限生成的 \(\mathbf{Z}\)-模，且 \(k''\) 是 \(k\) 的有限生成纯扩张。另一方面，由于 \(\mathbf{K}'\) 是 \(\mathbf{K}''\) 的有限代数扩张，\(\Gamma' / \Gamma\) 是有限群，且 \(k'\) 是 \(k''\) 的有限代数扩张（§ 8，第 1 节，引理 2）。这证明了推论。

推论 2. 令 \(h\) 和 \(h'\) 分别为 \(v\) 和 \(v'\) 的高。则

\[
s (v ^ {\prime} / v) + h ^ {\prime} \leqslant d (\mathrm{K} ^ {\prime} / \mathrm{K}) + h.\tag{10}
\]

由命题 3，\(h' \leqslant r(v'/v) + h\)。

推论 3. 设 \(\mathbf{K}'\) 是 \(\mathbf{K}\) 的有限生成扩张，\(\Gamma\) 同构于按字典序排列的 \(\mathbf{Z}^h\)，且公式 (10) 中取等号。则 \(\Gamma'\) 同构于按字典序排列的 \(\mathbf{Z}^{h'}\)，且 \(k'\) 是 \(k\) 的有限生成扩张。

若 (10) 中取等号，则 (9) 中也取等号，从而 \(k'\) 是 k 的有限生成扩张，且 \(\Gamma'\) 是有限生成 Z-模。此外，比较 (9) 与 (10)，可见 \(h' - h = r(\Gamma'/\Gamma)\)，从而 \(h' = r(\Gamma')\)，于是第 2 节的命题 4 表明 \(\Gamma'\) 同构于按字典序排列的 \(Z^{h'}\)。

推论 4. 设 \(v\) 是非真赋值（此时 \(k = K\)）。则

\[
h (\Gamma^ {\prime}) + d (k ^ {\prime} / \mathbf {K}) \leqslant r (\Gamma^ {\prime}) + \mathbf {D} (k ^ {\prime} / \mathbf {K}) \leqslant d (\mathbf {K} ^ {\prime} / \mathbf {K}).\tag{11}
\]

特别地，若 \(v'\) 的高为 1，则

\[
d (k ^ {\prime} / \mathbf {K}) \leqslant d (\mathbf {K} ^ {\prime} / \mathbf {K}) - 1;\tag{12}
\]

此外，若 \(\mathbf{K}'\) 是 \(\mathbf{K}\) 的有限生成扩张，且 (12) 中取等号，则 \(v'\) 是离散赋值，且 \(k'\) 是 \(\mathbf{K}\) 的有限生成扩张。

这是推论 1、2、3 的一系列特殊情形。


\exercisesection{1}

\begin{enumerate}
\def\labelenumi{\arabic{enumi}.}
\tightlist
\item
  设 K 为域；对于 K 的每个子环 A，令 L(A) 表示 A 的局部环 \(A_{p}\) 所成的集合，其中 p 遍历 A 的素理想（把这些局部环视为 K 的子环）。
\end{enumerate}

\begin{enumerate}
\def\labelenumi{(\alph{enumi})}
\item
  K 的局部子环 M 要支配 \(A_{p} \in L(A)\) 中的一个环，当且仅当 \(A \subset M\)；被 M 支配的局部环 \(A_{p}\) 是唯一的，并对应于素理想 \(p = m(M) \cap A\)。
\item
  设 M、N 是 K 的两个局部子环，P 是由 \(M \cup N\) 生成的 K 的子环。证明下列条件等价：
\end{enumerate}

\((\alpha)\) 存在素理想 \(\mathfrak{p}\)，属于 \(\mathbf{P}\)，使得 \(\mathfrak{m}(\mathbf{M}) = \mathfrak{p} \cap \mathbf{M}\)，\(\mathfrak{m}(\mathbf{N}) = \mathfrak{p} \cap \mathbf{N}\)。

(β) 理想 \(\mathfrak{a}\) 是 \(\mathbf{P}\) 中由 \(\mathfrak{m}(\mathbf{M}) \cup \mathfrak{m}(\mathbf{N})\) 生成的，且不同于 \(\mathbf{P}\)。

(γ) 存在 K 的一个同时支配 M 和 N 的局部子环 Q。

若这些条件满足，则称 M 和 N 相关联。

\begin{enumerate}
\def\labelenumi{(\alph{enumi})}
\setcounter{enumi}{2}
\tightlist
\item
  设 A、B 是 K 的两个子环，C 是由 A ∪ B 生成的 K 的子环。证明下列条件等价：
\end{enumerate}

\((\alpha)\) 对每个包含 A、B 的局部环 \(\mathbf{Q} \subset \mathbf{K}\)，都有 \(\mathbf{A}_{\mathfrak{p}} = \mathbf{B}_{\mathfrak{q}}\)，其中 \(\mathfrak{p} = \mathfrak{m}(\mathbf{Q}) \cap \mathbf{A}\)，\(\mathfrak{q} = \mathfrak{m}(\mathbf{Q}) \cap \mathbf{B}\)。

(β) 对 C 的每个素理想 r，都有 \(A_{p} = B_{q}\)，其中 \(p = r \cap A\)，\(q = r \cap B\)。

(γ) 若 M ∈ L(A) 且 N ∈ L(B) 相关联，则二者相同。

(δ) L(A) ∩ L(B) = L(C)。

（使用 (a) 和 (b)。）

\begin{enumerate}
\def\labelenumi{\arabic{enumi}.}
\setcounter{enumi}{1}
\item
  对于整环 A，A 是赋值环的充要条件是 A 的每个理想都是不可约的（《代数》第二章 § 2 习题 16）。
\item
  证明每个赋值环都是凝聚的（第一章 § 2 习题 12）。
\item
  设 K 为域，F 是 K 的赋值环的集合，按包含关系全序。证明 F 中各环的交是 K 的赋值环。
\item
  设 K 为域，A 是 K 的整闭子域，K 是其分式域，且 \((\mathbf{A}_{\alpha})\) 是 K 的一个赋值环族，其交为 A。若 L 是 K 的扩域，则 A 在 L 中的整闭包等于 L 中支配某个 \(A_{\alpha}\) 的赋值环之交。
\end{enumerate}

¶ 6. 设 R 为环，A 为 R 的子环，使 S = R - A 非空且 S 中两个元素之积属于 S。

\begin{enumerate}
\def\labelenumi{(\alph{enumi})}
\item
  设 \(a, a'\) 是 \(\mathbf{A}\) 中的两个元素，\(s, s'\) 是 \(S\) 中的两个元素，且 \(sa \in \mathbf{A}\)、\(s'a' \in \mathbf{A}\)。证明 \(sa', s'a\) 中至少有一个属于 \(\mathbf{A}\)。由此证明，集合 \(p_1\) 由 \(a \in \mathbf{A}\) 组成，其中存在 \(s \in S\) 使 \(sa \in \mathbf{A}\)，且它是 \(\mathbf{A}\) 的素理想。
\item
  证明集合 \(\mathfrak{p}_0\) 由 \(a \in \mathbf{A}\) 组成，其中 \(sa \in \mathbf{A}\) 对所有 \(s \in \mathbf{S}\) 成立，且它是 \(\mathbf{R}\) 和 \(\mathbf{A}\) 的素理想；它是 \(\mathbf{R}\) 的最大且包含于 \(\mathbf{A}\) 的理想，且 \(\mathfrak{p}_0 \subset \mathfrak{p}_1\)。
\item
  集合 \(\mathfrak{n}\) 由 \(x \in \mathbb{R}\) 组成，其中存在 \(s \in \mathbb{S}\) 使 \(sx = 0\)，既是 \(\mathbb{R}\) 的理想，也是 \(\mathbb{A}\) 的理想，且 \(\mathfrak{n} \subset \mathfrak{p}_0\)。
\item
  环 \(\mathbf{A}\) 在 \(\mathbb{R}\) 中整闭。
\item
  若 \(\mathbf{R}\) 是域，\(\mathbf{A}\) 是 \(\mathbf{R}\) 的赋值环，\(p_0 = (0)\)，且 \(p_1\) 是 \(\mathbf{A}\) 的极大理想。
\end{enumerate}

¶ 7. 设 R 为环。若有序对 (A, p) 由 R 的子环 A 及 A 的素理想 p 组成，且不存在由 R 的子环 A' 及 A' 的素理想 p' 组成的有序对满足 A ⊂ A'、A' ≠ A 且 p' ∩ A = p，则称 (A, p) 在 R 中是极大的。

\begin{enumerate}
\def\labelenumi{(\alph{enumi})}
\item
  对于 R 的每个子环 B 及 B 的每个素理想 q，存在在 R 中极大的有序对 \((A, p)\)，使 \(A \supset B\) 且 \(q = p \cap B\)。
\item
  有序对 (A, p) 在 R 中极大的充要条件是：对于所有 \(s \in R - A\)，存在元素 \(c_{i} \in p\) 的有限族（\(1 \leqslant i \leqslant n\)），使元素 \(b = c_{1}s + c_{2}s^{2} + \cdots + c_{n}s^{n}\) 属于 A - p（使用《代数》第二章 § 2 第 5 节命题 11 的推论 3）。
\item
  证明若有序对 \((\mathbf{A},\mathfrak{p})\) 在 \(\mathbf{R}\) 中极大，则 \(\mathbf{R} - \mathbf{A}\) 中两个元素之积属于 \(\mathbf{R} - \mathbf{A}\)。（若 \(s,s'\) 属于 \(\mathbf{R} - \mathbf{A}\) 且 \(ss' = a\in \mathbf{A}\)，考虑满足下述条件的最小整数 \(n,m\)：存在 \(c_{i}\in \mathfrak{p}\)（\(1\leqslant i\leqslant n\)）使 \(b = \sum_{i}c_{i}s^{i}\in \mathbf{A} - \mathfrak{p}\)，以及 \(c_{j}'\in \mathfrak{p}\)（\(1\leqslant j\leqslant m\)）使
\end{enumerate}

\[
b ^ {\prime} = \sum_ {j} c _ {j} ^ {\prime} s ^ {\prime j} \in A - p;
\]

  例如设 \(n \geqslant m\)，考察 \(bb' \in A - p\)，由整数 \(n\) 的定义得到矛盾。）

\begin{enumerate}
\def\labelenumi{(\alph{enumi})}
\setcounter{enumi}{3}
\item
  在 (c) 的假设下，令 \(S = R - A\)；证明理想 \(p_1\) 是 \(A\) 中由 \(a \in A\) 构成的集合，其中存在 \(s \in S\) 使 \(sa \in A\)，并且它包含于 \(p\)（使用 (b)）。同理，理想 \(p_0\) 是 \(A\) 中由 \(a \in A\) 构成的集合，其中 \(sa \in A\) 对所有 \(s \in S\) 成立，并且它包含于 \(p\)（习题 6 (b)）。记 \(A' = A / p_0\)、\(p' = p / p_0\) 和 \(R' = R / p_0\)；有序对 \((A', p')\) 于是是整环 \(\mathbf{R}'\) 中的极大有序对，并且对于每个非零 \(a' \in \mathbf{A}'\)，存在 \(s' \in \mathbf{S}' = \mathbf{R}' - \mathbf{A}'\)，使 \(s'a' \in \mathbf{S}'\)。
\item
  设 \((\mathbf{A},\mathfrak{p})\) 为整域 \(\mathbf{R}\) 中的极大有序对，使得对于每个非零 \(a\in \mathbf{A}\)，存在 \(s\in S = R - A\) 使 \(sa\in S\)。记 \(\mathrm{S}_0 = \mathrm{S}\cup \{1\}\)，并令 \(\mathbb{R}_0\) 表示分式环 \(\mathrm{S}_0^{-1}\mathrm{R}\)；典范同态 \(\mathbb{R}\to \mathbb{R}_0\) 是单射，因而将 \(\mathbb{R}\) 识别为 \(\mathbb{R}_0\) 的子环。证明 \(\mathbb{R}_0\) 是域，并可识别为 \(\mathbb{R}\) 的分式域。令 \((\mathbf{B},\mathfrak{q})\) 为 \(\mathbb{R}_0\) 中的极大有序对，使 \(\mathbf{B}\supset \mathbf{A}\) 且 \(\mathfrak{q}\cap \mathbf{A} = \mathfrak{p}\)。证明 \(\mathbf{B}\cap \mathbf{R} = \mathbf{A}\)；换言之，\(\mathbf{A}\) 是 \(\mathbf{R}\) 与 \(\mathbb{R}_0\) 中某个赋值环的交，且 \(\mathfrak{p}\) 是 \(\mathbf{A}\) 与该赋值环的极大理想的交。
\item
\end{enumerate}

¶ 8. 给定环 A，若系数属于 A 的多项式 P(X \(_{1}\) , \ldots, X \(_{n}\) ) 具有形式 X \(^{\alpha}\) + ∑ \(_{\beta} c_{\beta}\) X \(^{\beta}\)，则称其为主导的，其中 β \textless{} α 属于乘积有序集 N \(^{n}\)，并且对所有满足 c \(_{\beta}\) ≠ 0 的 β 都如此。在环 R 中，若子环 A 对每个主导多项式 P ∈ A{[}X \(_{1}\) , \ldots, X \(_{n}\) {]}（任意 n）都满足：P(s \(_{1}\) , \ldots, s \(_{n}\) ) ≠ 0 对所有元素 s \(_{i}\) ∈ R − A（1 ≤ i ≤ n），则称 A 为旁赋值的。

\begin{enumerate}
\def\labelenumi{(\alph{enumi})}
\item
  证明若 \(\mathbf{A}\) 在 \(\mathbb{R}\) 中是旁赋值的，则 \(\mathbf{S} = \mathbf{R} - \mathbf{A}\) 中两个元素之积属于 \(\mathbf{S}\)。
\item
  设 A 是环 R 的子环，且 S = R - A 中两个元素之积属于 S；令 \(p_{0}\) 为 R 和 A 中由 \(a \in A\) 组成的素理想，其中 \(sa \in A\) 对所有 \(s \in S\) 成立（习题 6 (b)）。A 在 R 中为旁赋值的充要条件是 \(A/p_{0}\) 在 \(R/p_{0}\) 中为旁赋值的。
\item
  设 R 是环，\(h: \mathbf{R} \to \mathbf{R}'\) 是环同态。若 \(\mathbf{A}'\) 是 \(\mathbf{R}'\) 的旁赋值子环，\(\mathbf{A} = \bar{h}^{-1}(\mathbf{A}')\) 在 R 中旁赋值。由此证明：若 \((\mathbf{A}, p)\) 是 R 中的极大有序对（习题 7），则 A 在 R 中旁赋值（使用 (b) 及习题 7 (d)、(e)）。
\item
  设 \(\mathbf{R}'\) 为环，\(\mathbf{R}\) 为 \(\mathbf{R}'\) 的子环。若 \(\mathbf{A}\) 是 \(\mathbf{R}\) 的旁赋值子环，证明存在旁赋值子环 \(\mathbf{A}'\)，它属于 \(\mathbf{R}'\)，使 \(\mathbf{A}' \cap \mathbf{R} = \mathbf{A}\)。（若 \(S = R - A\) 且 \(S_0 = S \cup \{1\}\)，考察环 \(T' = S_0^{-1}R'\) 及典范同态 \(h: R' \to T'\)。令 \(\mathbf{B}\) 为 \(T'\) 中由 \(h(A) \cup (h(S))^{-1}\) 生成的子环，令 \(\mathfrak{b}\) 为 \(\mathbf{B}\) 中由 \((h(S))^{-1}\) 生成的理想。证明 \(\mathfrak{b} \neq \mathbf{B}\)；取极大有序对 \((\mathbf{A}^{\prime \prime}, \mathfrak{p}^{\prime \prime})\)，它属于 \(\mathbf{T}'\)，使 \(\mathbf{B} \subset \mathbf{A}''\) 且 \(\mathfrak{b} \subset \mathfrak{p}''\)，并令 \(\mathbf{A}' = \overline{h}^{-1}(\mathbf{A}'')\)。
\end{enumerate}

由此证明，若 R 是整环，则 R 的旁赋值子环正好是 R 与其分式域的赋值环的交。

¶ 9. (a) 设 R 为环，A 为 R 的子环，x 为 R 中元素，S 为 R 中由 \(x^{k}\)（\(k \geqslant 0\)）组成的集合。证明 x 在 A 上整当且仅当 \(x/1 \in A[1/x]\)（在环 \(S^{-1}R\) 中）。若 x 不在 A 上整，证明 \(A[1/x]\) 有一个包含 1/x 的极大理想 m，且 m 在典范同态 \(A \to A[1/x]\) 下的逆像是 A 的极大理想。

\begin{enumerate}
\def\labelenumi{(\alph{enumi})}
\setcounter{enumi}{1}
\tightlist
\item
  证明 A 在 R 中的整闭包 A' 等于包含 A 的 R 的旁赋值子环（习题 8）之交。（为证明 A' 包含于这些环中的每一个，使用习题 8 (a) 和习题 6 (d)；为证明若 \(x \in R\) 不在 A 上整，则存在 R 的一个在 R 中旁赋值的子环 B 使 \(x \notin B\)，使用 (a)、习题 7 (a) 和习题 8 (c)，仿照第 3 节定理 3 进行论证。）
\end{enumerate}

\exercisesection{2}

\begin{enumerate}
\def\labelenumi{\arabic{enumi}.}
\item
  设 \(\Omega\) 是域 K 的一个扩域，E、F 是包含于 \(\Omega\) 中的 K 的两个扩域，且在 K 上线性不交，L 是包含于 E 中的 K 的扩域。令 \(f\) 为 E 到 L 的位，它在 K 上化为恒等自同构。证明存在唯一的位，定义在 K(E \(\cup\) F) 上，取值于 K(L \(\cup\) F)，延拓 \(f\)，并在 F 上化为恒等自同构。（若 A 是 \(f\) 的环，注意 \(f\) 可唯一延拓为同态 \(g\)，从 F{[}A{]} 到 K(L \(\cup\) F)，在 F 上化为恒等，并且若 m 是 \(f\) 的理想，则每个满足 \(z \in\) F{[}A{]} 且 \(g(z) = 0\) 的元素 z 都可写成 \(xu\)，其中 \(x \in m\)，\(u \in\) F{[}A{]} 且 \(g(u) \neq 0\)。最后注意 K(E \(\cup\) F) 是 F{[}A{]} 的分式域。）
\item
  设 K、K' 是域 k 的两个扩域，f 和 \(f'\) 分别是 K 和 \(K'\) 到同一代数闭域 L 的位。假设 f 与 \(f'\) 在 k 上的限制相同。证明存在 K 与 \(K'\) 的复合扩域 (F, i, i')（《代数》第八章 § 8）以及 F 到 L 的位 g，使得 \(f(x) = g(i(x))\) 对 \(x \in \mathbf{K}\) 成立，且 \(f'(x) = g(i'(x'))\) 对 \(x' \in \mathbf{K}'\) 成立。（令 V、\(V'\) 分别为 f、\(f'\) 的环，A 为它们与 k 的公共交，且 h: \(V \otimes_{A} V' \to L\) 为满足 \(h(a \otimes a') = f(a)f'(a')\)（其中 \(a \in V\)、\(a' \in V'\)）的同态。利用 V 和 \(V'\) 是平坦 A-模这一事实（§ 3 第 5 节引理 1），证明若 \(a \neq 0\)、\(a' \neq 0\)，则 \(a \otimes a' = (a \otimes 1)(1 \otimes a')\) 不是零因子；若 S 是 B = V \(\otimes_{A} V'\) 中由这些元素组成的乘法子集，
\end{enumerate}

\(S^{-1}B\) 含有子域 \(K_{1}, K_{1}'\)，它们分别同构于 K 和 \(K'\)，并且它是由 \(K_{1} \cup K_{1}'\) 生成的环。证明若 q 是 h 的核，则 B 中存在素理想 p，使 \(p \subset q\) 且 \(p \cap S = \varnothing\)，并证明 F 可取为 \(S^{-1}B/S^{-1}p\) 的分式域。

\begin{enumerate}
\def\labelenumi{\arabic{enumi}.}
\setcounter{enumi}{2}
\tightlist
\item
  \begin{enumerate}
  \def\labelenumii{(\alph{enumii})}
  \tightlist
  \item
    设 \(\mathbf{K}\) 为域，\(f\) 是 \(\mathbf{K}\) 上取值于可序域 \(\mathbf{L}\) 的位（《代数》第六章 § 2 习题 8）；证明 \(\mathbf{K}\) 可赋序。（注意，对于元素有限族 \((a_i)_{1\leq i\leq n}\) 属于 \(\mathbf{K}\)，总存在指标 \(j\)，使 \(f(a_i / a_j)\neq \infty\) 对所有 \(i\) 成立。）
  \end{enumerate}
\end{enumerate}

若 \((x_{i})_{1\leqslant i\leqslant n}\) 是 K 中元素的序列，且 \(f\left(\sum_{i = 1}^{n}x_{i}^{2}\right)\neq \infty\)，证明 \(f(x_{i})\neq \infty\) 对所有 \(i\) 成立。

\begin{enumerate}
\def\labelenumi{(\alph{enumi})}
\setcounter{enumi}{1}
\item
  设 K 为有序域，G 为 K 的子群；证明环 \(F(G)\)（由 K 中相对于 G 不是无限大的元素组成；同上，习题 11）是 K 的赋值环，而理想 \(I(G)\)（由 \(F(G)\) 中相对于 G 无穷小的元素组成）是此环的极大理想；因而 K 上存在一个取值于 \(k(G) = F(G)/I(G)\) 的位，它在 G 上化为恒等；称此位为 K 关于 G 的典范位。
\item
  证明若 K 中不存在异于 G 且与 G 可比较的 G 的扩域，则 \(k(G)\) 在 G 上代数（注意，若 \(t \in K\) 不属于 G，则域 \(G(t)\) 含有相对于 G 无穷小的元素 \(u \neq 0\)，且 \(G(t)\) 是 \(G(u)\) 的代数扩张）。
\end{enumerate}

¶ 4. 若有序域 K 对于 K 中所有 \(x \geqslant 0\) 都存在 \(y \in \mathbf{K}\) 使 \(x = y^2\)，则称 K 为欧氏的。

\begin{enumerate}
\def\labelenumi{(\alph{enumi})}
\item
  设 K 为欧氏有序域，G 为 K 的子域，且 K 中不存在异于 G 且与 G 可比较的 G 的扩域。证明：若 f 是 K 到有序域 L 的位，L 是 G 的代数扩域，且 f 在 G 上化为恒等，则 f 等价于 K 关于 G 的典范位。（可设 f 取值于 L 的一个极大有序代数扩域 N。注意 G 是欧氏的，并且若 \(x \in G\)、\(y \in K\) 且 x \textless{} y，则 \(f(y) = \infty\) 或 \(f(y) > x\)。若 A 和 p 是 f 的环和理想，依次证明 \(p \subset I(G)\)、\(F(G) \subset A\)、\(I(G) \subset p\) 以及 \(A \subset F(G)\)，并注意 N 与 G 可比较。）
\item
  设 K 为欧氏有序域，f 为一个取值于极大有序域 L 的 K 上的位，A 和 p 为 f 的环和理想。~令 G 为满足 \(E \cap p = 0\) 的 K 的子域 E 中的极大子域。证明 \(G \subset A\)，且 G 在 K 中代数闭，因而是欧氏的。证明 K 中不存在异于 G 且与 G 可比较的 G 的扩域。此外，L 中由 \(f(A)\) 生成的子域在 \(f(G)\) 上代数，因而 f 等价于 K 关于 G 的典范位。
\item
  ¶ 6. 设 K 为极大有序域，L 为 K 的异于 K 的极大有序子域，且 K 与 L 可比较（《代数》第六章 § 2 习题 15）。证明不存在取值于 L、在 L 上化为恒等自同构的 K 的位 f。
\end{enumerate}

¶ 5. 设 K 为域，f 为一个取值于极大有序域 L 的 K 上的位。

\begin{enumerate}
\def\labelenumi{(\alph{enumi})}
\item
  对所有 \(\alpha \in \mathbf{K}\)，证明 \(f\) 存在一个延拓，要么延拓到 \(\mathbf{K}(\sqrt{\alpha})\)，要么延拓到 \(\mathbf{K}(\sqrt{-\alpha})\)，并且该延拓是取值于 \(\mathbf{L}\) 的位。（根据以下两种情形之一进行讨论：存在 \(a \in \mathbf{K}\) 使 \(f(a^2\alpha)\) 既非 0 也非 \(\infty\)，或者 \(f(x^2\alpha)\) 等于 0 或 \(\infty\) 对所有 \(x \in \mathbf{K}\) 成立。）
\item
由 (a) 推出，存在 K 的一个扩域 E，它是极大有序域，并且 f 到 E 的延拓是取值于 L 的 E 上的位。（将其归结为 K 为欧氏的情形，并使用习题 4 (a)。）
\end{enumerate}

\begin{enumerate}
\def\labelenumi{\arabic{enumi}.}
\setcounter{enumi}{5}
\tightlist
\item
  设 A 为整闭整环，K 为其分式域，p 为 A 的素理想，k 为 A/p 的分式域。~证明若 \(A_{p}\) 不是赋值环，则 K 中存在支配 \(A_{p}\) 的赋值环 V，且其剩余域是 k 的超越扩张（使用《代数》第五章 § 2 习题 14(b)）。
\end{enumerate}

\exercisesection{3}

¶ 1. (a) 设 L 为域，T 为不定元；在形式幂级数域 L((T)) 中，考察级数

\[
s = c _ {0} + c _ {1} \mathrm{T} ^ {1!} + c _ {2} \mathrm{T} ^ {2!} + \dots + c _ {n} \mathrm{T} ^ {n!} + \dots
\]

其中 \(c_{i} \in L\) 均 \(\neq 0\)。证明 s 在 \(\mathrm{L}(\mathrm{T})\) 上超越；\(\mathrm{L}((\mathrm{T}))\) 是其扩域。（反证，假设存在关系 \(a_{0}(\mathrm{T}) + a_{1}(\mathrm{T})s + \cdots + a_{q}(\mathrm{T})s^{q} = 0\)，其中 \(a_{i}(\mathrm{T})\) 是 \(\mathrm{L}(\mathrm{T})\) 中的多项式且 \(a_{q} \neq 0\)；构造下述元素所满足的方程

\[
s ^ {\prime} = s - c _ {0} - c _ {1} \mathrm{T} ^ {1!} - \dots - c _ {n - 1} \mathrm{T} ^ {(n - 1)!}
\]

并证明当 \(n\) 足够大时，其常数项是次数 \(< n!\) 的多项式，从而矛盾。）

\begin{enumerate}
\def\labelenumi{(\alph{enumi})}
\setcounter{enumi}{1}
\tightlist
\item
  设 \(k\) 为域，\(\mathbf{X}, \mathbf{Y}, \mathbf{Z}\)、\(\mathrm{T}\) 为四个不定元，\(\mathbf{K} = k(\mathbf{X}, \mathbf{Y}, \mathbf{Z})\)，\(\mathbf{E}\) 为 \(k(\mathbf{X})\) 的代数闭包；在 \(\mathbf{E}\) 中，对所有 \(n\)，令 \(x_{n}\) 表示满足 \(x_{n}^{n} = \mathbf{X}\) 的元素。在形式幂级数域 \(\mathrm{E}((\mathrm{T}))\) 中，考察元素
\end{enumerate}

\[
s = x _ {1} \mathrm{T} + x _ {2} \mathrm{T} ^ {2!} + \dots + x _ {n} \mathrm{T} ^ {n!} + \dots
\]

由 (a)，元素 \(s\) 和 \(\mathbf{T}\) 在 \(\mathbf{E}\) 上代数独立。考察同态 \(f\colon \mathbf{K} \to \mathbf{E}((\mathbf{T}))\)，满足 \(f(\mathbf{X}) = \mathbf{X}, f(\mathbf{Y}) = \mathbf{T}, f(\mathbf{Z}) = s\)。若 \(v\) 是 \(\mathbf{E}((\mathbf{T}))\) 上由形式幂级数的阶定义的离散赋值（第 3 节例 3），则 \(v \circ f\) 是 \(\mathbf{K}\) 上的离散赋值；证明 \(v \circ f\) 的剩余域是子域 \(k(x_1, x_2, \ldots, x_n, \ldots)\)，它属于 \(\mathbf{E}\)，因而在 \(k\) 上是无限的。

\begin{enumerate}
\def\labelenumi{\arabic{enumi}.}
\setcounter{enumi}{1}
\tightlist
\item
  设 \(\Gamma\) 为全序群，\(k\) 为域。
\end{enumerate}

\begin{enumerate}
\def\labelenumi{(\alph{enumi})}
\item
  设 A、B 是 \(\Gamma\) 的两个子集，按诱导序良序；证明集合 \(A + B\) 是良序的，并且对于每个 \(\gamma \in A + B\)，满足 \((\alpha, \beta) \in A \times B\) 且 \(\alpha + \beta = \gamma\) 的有序对的集合是有限的。（为证明第一个断言，反证：否则可以定义严格递减序列 \((\gamma_n)\)，其各项属于 \(A + B\)，使 \(\gamma_n = \alpha_n + \beta_n\)，其中 \(\alpha_n \in A\)、\(\beta_n \in B\)，序列 \((\alpha_n)\) 严格递增而序列 \((\beta_n)\) 严格递减；使用《集合论》第三章 § 4 习题 3 得出结论。）
\item
  设 \(S(\Gamma, k)\) 为下述族的集合：\(x = (x_{\alpha})_{\alpha \in \Gamma}\)，满足 \(x_{\alpha} \in k\) 对所有 \(\alpha \in \Gamma\) 成立，且集合 \(\alpha\)（满足 \(x_{\alpha} \neq 0\)）是 \(\Gamma\) 的良序子集。证明 \(S(\Gamma, k)\) 是 \(k^{\Gamma}\) 的加法子群；此外，若 \(x = (x_{\alpha}), y = (y_{\alpha})\) 是 \(S(\Gamma, k)\) 中的两个元素，则集合 \(C\) 由 \(\alpha + \beta\) 组成，其中有序对 \((\alpha, \beta)\) 满足 \(x_{\alpha} \neq 0\) 且 \(y_{\beta} \neq 0\)；由 (a)，它是良序的，并且对每个 \(\gamma \in C\)，元素 \(z_{\gamma} = \sum_{\alpha + \beta = \gamma} x_{\alpha} y_{\beta}\) 在 \(k\) 中有定义；若置 \(z_{\alpha} = 0\) 对 \(\alpha \notin C\)，则元素 \(z = (z_{\alpha}) \in k^{\Gamma}\) 属于 \(S(\Gamma, k)\)；将其记为 \(xy\)。证明在乘积群 \(k^{\Gamma}\) 上赋予上述乘法和加法后，集合 \(S(\Gamma, k)\) 是域；此外，对于 \(x = (x_{\alpha}) \neq 0\) 属于 \(S(\Gamma, k)\) 的任意元素，令 \(\lambda\) 为 \(\alpha \in \Gamma\) 中使 \(x_{\alpha} \neq 0\) 的最小者；若记 \(v(x) = \lambda\)（并令 \(v(0) = +\infty\)），证明 \(v\) 是 \(S(\Gamma, k)\) 上的赋值，剩余域为 \(k\)，值群为 \(\Gamma\)。称 \(v\) 为域 \(S(\Gamma, k)\) 上的典范赋值。域 \(k\) 典范地识别为 \(S(\Gamma, k)\) 的一个子域。对所有 \(\alpha \in \Gamma\)，令 \(u_{\alpha}\) 表示元素 \((x_{\lambda})_{\lambda \in \Gamma}\)，满足 \(x_{\alpha} = 1, x_{\lambda} = 0\) 对 \(\lambda \neq \alpha\)；于是，对每个非零 \(z \in S(\Gamma, k)\)，存在唯一的有序对 \((t, \alpha) \in k \times \Gamma\)，使 \(v(z) = \alpha\) 且 \(v(z - tu_{\alpha}) > \alpha\)。称 \(tu_{\alpha}\) 为 \(z\) 的主导项。
\end{enumerate}

¶ 3. 设 K 为域，w 为 K 上的赋值，Γ 为 w 的值群，k 为 w 的剩余域。构造相应的域 S(Γ, k)，并沿用习题 2 的记号。对于所有 \(t \in k\)，令 \(t^{0}\) 为 w 的赋值环中类 t 的一个元素；对于所有 \(\alpha \in \Gamma\)，令 \(u_{\alpha}^{0}\) 为 K 中满足 \(w(u_{\alpha}^{0}) = \alpha\) 的元素。

\begin{enumerate}
\def\labelenumi{(\alph{enumi})}
\tightlist
\item
  设 M 是 S(Γ, k) 的一个子集，包含元素 \(tu_{\alpha}\)，其中每个有序对 \((t, \alpha) \in k \times \Gamma\) 都对应一个这样的元素；假设存在双射 \(x \mapsto x^{0}\)，它把 M 映到 K 的一个子集 \(M^{0}\)，且满足下列条件：(1) \((tu_{\alpha})^{0} = t^{0}u_{\alpha}^{0}\) 对每个有序对 \((t, \alpha) \in k \times \Gamma\) 成立；(2) 若 \(x = (x_{\alpha})\) 属于 M，且 \(C_{x}\) 是 Γ 的良序子集，由满足 \(x_{\alpha} \neq 0\) 的 α 组成，则对 \(C_{x}\) 的每个区段，元素 \((x_{\lambda})_{\lambda \in D}\) 属于 M；(3) 若 x、y 是 M 的两个元素，且 \(tu_{\alpha}\) 是 x - y 的主导项，则 \(w(x^{0} - y^{0} - t^{0}u_{\alpha}^{0}) > w(x^{0} - y^{0})\)。若 \(z'\) 是 K 中不属于 \(M^{0}\) 的元素，证明存在 S(Γ, k) 中的元素 \(z \notin M\)，使得在 M 上与 \(x \mapsto x^{0}\) 重合且把 \(z'\) 映到 z 的映射也满足上述条件。（注意，若 \(w(z') = \alpha\)，存在 \(t^{0}\)，使
\end{enumerate}

\[
w (z ^ {\prime} - t ^ {0} u _ {\alpha} ^ {0}) > \alpha ;
\]

令 \(z_{\alpha} = tu_{\alpha}\)，然后证明通过超限归纳可以确定 \(z_{\beta}\)（其中 \(\beta > \alpha\)），使 \(z = (z_{\lambda})_{\lambda \in \Gamma}\) 解决此问题。）

\begin{enumerate}
\def\labelenumi{(\alph{enumi})}
\setcounter{enumi}{1}
\tightlist
\item
  由 (a) 推出 \(\operatorname{Card}(\mathbf{K}) \leqslant \operatorname{Card} S(\Gamma, k)\)。
\end{enumerate}

\begin{enumerate}
\def\labelenumi{\arabic{enumi}.}
\setcounter{enumi}{3}
\tightlist
\item
  设 A 是局部整环，K 是其分式域，m 是其极大理想，v 是 K 上的赋值，且 v 的环 B 支配 A。假设 m 是有限生成理想，或者 v 是离散赋值；则 m 中存在元素 x，使 \(v(x) = \inf_{y \in \mathfrak{m}} v(y)\)。令 \(A'\) 为 K 中由
\end{enumerate}

A 以及 \(yx^{-1}\)（其中 y 遍历 m）生成的子环；令 \(A_{1}\) 为 \(A'\) 的局部环，相对于素理想 \(p \cap A'\)，其中 p 是赋值 v 的理想。证明环 \(A_{1}\) 不依赖于满足上述条件的元素 x 的选取（\(A_{1}\) 称为“在 \(v'\) 方向上 A 的素二次变换”）；若 A 是 Noether 环，则 \(A_{1}\) 也是 Noether 环。

¶ 5. 设 k 为域，K = k(X, Y) 是 k 上两个不定元的有理函数域。记 \(x_{1} = X\)、\(y_{1} = Y\)，并对 \(n \geqslant 1\) 归纳定义 K 中的元素 \(x_{n+1} = y_{n}\) 和 \(y_{n+1} = x_{n}y_{n}^{-1}\)。

\begin{enumerate}
\def\labelenumi{(\alph{enumi})}
\tightlist
\item
  写 \(\mathbf{A}_n = k[x_n, y_n]\)；序列 \((\mathbf{A}_n)\) 递增，并且若
\end{enumerate}

\[
\mathfrak {m} _ {n} = \mathrm{A} _ {n} x _ {n} + \mathrm{A} _ {n} y _ {n},
\]

\(m_{n}\) 是 \(A_{n}\) 的极大理想，且 \(m_{n} = A_{n} \cap m_{n+1}\)；在环 \(A = \bigcap_{n} A_{n}\) 中，\(m = \bigcap_{n} m_{n}\) 是极大理想。

\begin{enumerate}
\def\labelenumi{(\alph{enumi})}
\setcounter{enumi}{1}
\item
  设 \(v\) 为 \(\mathbf{K}\) 上由 \(v(\mathbf{X}) = \rho = (1 + \sqrt{5}) / 2\)、\(v(\mathbf{Y}) = 1\) 定义的赋值；令 \(\mathbf{B}\) 和 \(\mathfrak{p}\) 为 \(v\) 的环和理想；证明 \(\mathfrak{m} = \mathfrak{p} \cap \mathbf{A}\)（注意 \(\rho^2 - \rho - 1 = 0\)，并用此计算 \(v(y_n)\)）。证明对于单项式 \(\mathbf{X}^i\mathbf{Y}^j\)（\(i, j\) 为正整数或负整数）属于 \(\mathbf{A}\) 的充要条件是 \(i\left(\frac{1 + \sqrt{5}}{2}\right) + j > 0\)；因此，对于每个单项式 \(z = \mathbf{X}^i\mathbf{Y}^j\)，要么 \(z \in \mathbf{A}\)，要么 \(1/z \in \mathbf{A}\)。
\item
  证明 \(B = A_{m}\)（将 \(t \in B\) 写成 X、Y 中两个多项式之商的形式，并证明在这两个多项式中，v 取最小值的单项式……）。
\item
  证明对所有 \(n\)，\((\mathbf{A}_{n+1})_{\mathfrak{m}_{n+1}}\) 是 \((\mathbf{A}_n)_{\mathfrak{m}_n}\) 在 \(v\) 方向上的素二次变换（练习 4）。
\end{enumerate}

\begin{enumerate}
\def\labelenumi{\arabic{enumi}.}
\setcounter{enumi}{5}
\tightlist
\item
  \begin{enumerate}
  \def\labelenumii{(\alph{enumii})}
  \tightlist
  \item
    设 \(\mathbf{A}\) 为赋值环。将《代数》第七章 § 4 第 1 至 4 节的结果（不变因子理论）推广到有限生成 A-模。
  \end{enumerate}
\end{enumerate}

\begin{enumerate}
\def\labelenumi{(\alph{enumi})}
\setcounter{enumi}{1}
\item
  设 K 为域，v 为 K 上的赋值，A 为 v 的环，且 \(\Gamma\) 为 v 的阶群。令 \(U = (\alpha_{ij})\) 为 A 中元素构成的 n 阶方阵，且其行列式为 0；证明存在 \(\lambda \in \Gamma\)，使得对于每个由 A 中元素构成的 n 阶方阵 \(U' = (\alpha_{ij}')\)，若 \(v(\alpha_{ij}' - \alpha_{ij}) > \lambda\) 对每个有序对 \((i,j)\) 都成立，则 \(U'\) 的不变因子与 U 的相同。
\item
  在 (b) 的假设下，推广《代数》第九章 § 5 第 1 节的定理和习题 1 以及 § 6 习题 10 的下列结果（对于最后一个结果，假设 K 的特征不为 2 且 \(v(2) = 0\)）。
\end{enumerate}

\begin{enumerate}
\def\labelenumi{\arabic{enumi}.}
\setcounter{enumi}{6}
\item
  证明第二章 § 3 习题 2 (b) 中定义的整环 B 是局部环，其极大理想是主理想，但 B 不是赋值环。
\item
  证明若赋值环 A 是强 Lasker 环（第四章 § 2 习题 28），则 A 是域或离散赋值环（使用第四章 § 2 习题 29）。
\end{enumerate}
